% =====================================================================
%  TP Git & GitHub pour débutants -- BTS CIEL 1re année -- Gestion de projet
%
%  Un seul source, deux PDF (voir Makefile) :
%    Version élève   : pdflatex -jobname=TP_Git_debutant tp_git.tex
%    Version corrigé : pdflatex -jobname=TP_Git_debutant_corrige "\def\CORRIGE{}% =====================================================================
%  TP Git & GitHub pour débutants -- BTS CIEL 1re année -- Gestion de projet
%
%  Un seul source, deux PDF (voir Makefile) :
%    Version élève   : pdflatex -jobname=TP_Git_debutant tp_git.tex
%    Version corrigé : pdflatex -jobname=TP_Git_debutant_corrige "\def\CORRIGE{}% =====================================================================
%  TP Git & GitHub pour débutants -- BTS CIEL 1re année -- Gestion de projet
%
%  Un seul source, deux PDF (voir Makefile) :
%    Version élève   : pdflatex -jobname=TP_Git_debutant tp_git.tex
%    Version corrigé : pdflatex -jobname=TP_Git_debutant_corrige "\def\CORRIGE{}% =====================================================================
%  TP Git & GitHub pour débutants -- BTS CIEL 1re année -- Gestion de projet
%
%  Un seul source, deux PDF (voir Makefile) :
%    Version élève   : pdflatex -jobname=TP_Git_debutant tp_git.tex
%    Version corrigé : pdflatex -jobname=TP_Git_debutant_corrige "\def\CORRIGE{}\input{tp_git.tex}"
%  Compiler deux fois (numéro de dernière page dans le pied de page).
% =====================================================================
\documentclass[11pt,a4paper]{article}

\usepackage[utf8]{inputenc}
\usepackage[T1]{fontenc}
\usepackage[french]{babel}
\usepackage{lmodern}
\usepackage{microtype}
% Pas de ligatures -- et << >> dans les commandes (police à chasse fixe)
\DisableLigatures[-,<,>]{encoding=T1, family=tt*}
\usepackage{needspace}
\usepackage[margin=2.2cm,headheight=14pt]{geometry}
\usepackage{xcolor}
\usepackage{upquote}
\usepackage{listings}
\usepackage[most]{tcolorbox}
\usepackage{tikz}
\usetikzlibrary{positioning,arrows.meta,babel}
\usepackage{enumitem}
\usepackage{tabularx}
\usepackage{booktabs}
\usepackage{amssymb}
\usepackage{fancyhdr}
\usepackage{lastpage}
\usepackage{comment}
\usepackage[hidelinks]{hyperref}

% ---------------------------------------------------------------------
% Version élève / version corrigée
% ---------------------------------------------------------------------
\newif\ifcorrige
\ifdefined\CORRIGE\corrigetrue\fi

\definecolor{ciel}{RGB}{0,84,147}

\ifcorrige
  \newtcolorbox{correction}{breakable, colback=green!4, colframe=green!45!black,
    fonttitle=\bfseries\small, title=Correction, before skip=2mm, after skip=3mm}
  \newcommand{\reponse}[1]{}
\else
  \excludecomment{correction}
  \newcommand{\reponse}[1]{\par\smallskip{\color{gray!70}%
    \foreach \n in {1,...,#1}{\noindent\mbox{}\dotfill\par\vspace{2.5mm}}}}
\fi

% ---------------------------------------------------------------------
% Mise en forme
% ---------------------------------------------------------------------
\setlength{\parindent}{0pt}
\setlength{\parskip}{4pt}
\setlist{itemsep=2pt, topsep=3pt}

\pagestyle{fancy}
\fancyhf{}
\lhead{\small BTS CIEL 1 -- Gestion de projet}
\rhead{\small TP Git \& GitHub\ifcorrige{} -- \textcolor{red!70!black}{corrigé}\fi}
\cfoot{\small \thepage/\pageref{LastPage}}

% Commande en ligne : \detokenize permet d'écrire _ ~ : sans échappement
% (pas d'accents ni de # dans l'argument).
\newcommand{\cmd}[1]{\texttt{\textcolor{ciel!70!black}{\detokenize{#1}}}}
% Élément de l'interface web de GitHub
\newcommand{\gh}[1]{\textsf{\textbf{#1}}}

\lstdefinestyle{bash}{
  basicstyle=\ttfamily\small, backgroundcolor=\color{gray!8},
  frame=single, rulecolor=\color{gray!40}, framesep=4pt,
  xleftmargin=4pt, xrightmargin=4pt,
  columns=fullflexible, keepspaces=true, breaklines=true,
  showstringspaces=false, upquote=true, aboveskip=6pt, belowskip=6pt,
  literate={é}{{\'e}}1 {è}{{\`e}}1 {ê}{{\^e}}1 {à}{{\`a}}1 {â}{{\^a}}1
           {ù}{{\`u}}1 {û}{{\^u}}1 {ô}{{\^o}}1 {î}{{\^i}}1 {ç}{{\c{c}}}1
           {É}{{\'E}}1 {È}{{\`E}}1 {À}{{\`A}}1
}
\lstdefinestyle{sortie}{style=bash, backgroundcolor=\color{white},
  frame=leftline, framerule=2pt, rulecolor=\color{ciel!50}, basicstyle=\ttfamily\small\color{black!75}}
\lstset{style=bash}

\newtcolorbox{info}[1][À retenir]{breakable, colback=ciel!5, colframe=ciel!80!black,
  fonttitle=\bfseries, title=#1}
\newtcolorbox{attention}[1][Attention]{breakable, colback=red!4, colframe=red!60!black,
  fonttitle=\bfseries, title=#1}
\newtcolorbox[auto counter]{exercice}[1]{breakable, colback=white, colframe=orange!80!black,
  fonttitle=\bfseries, title={Exercice~\thetcbcounter~--~#1}}

\tikzset{
  zone/.style={draw, thick, rounded corners, minimum width=2.9cm, minimum height=1.4cm,
               align=center, fill=#1, font=\small},
  >={Stealth[length=2.5mm]}
}

\begin{document}

% =====================================================================
\begin{tcolorbox}[colback=ciel!6, colframe=ciel, arc=2mm]
  {\LARGE\bfseries TP -- Premiers pas avec Git et GitHub}\\[1mm]
  {\large Sauvegarder et partager son travail, la bonne routine}\\[2mm]
  BTS CIEL -- 1\textsuperscript{re} année \hfill Gestion de projet
  \ifcorrige\\[2mm]{\bfseries\color{red!70!black} VERSION CORRIGÉE -- DOCUMENT ENSEIGNANT}\fi
  \tcblower
  \begin{tabularx}{\linewidth}{@{}>{\bfseries}l X@{}}
    Durée      & 2 h à 3 h \\
    Prérequis  & ouvrir un terminal, utiliser \cmd{cd} et \cmd{ls}, éditer un fichier avec \cmd{nano} ou VS Code \\
    Matériel   & poste avec Git (Git Bash sous Windows), navigateur web, adresse e-mail pour créer un compte GitHub \\
  \end{tabularx}
\end{tcolorbox}

\ifcorrige\else
\vspace{2mm}
\textbf{Nom :} \makebox[6cm]{\dotfill} \hfill \textbf{Prénom :} \makebox[6cm]{\dotfill}
\vspace{2mm}
\fi

\textbf{Objectif.} Ce TP ne fait pas de vous des experts de Git. Il vous apprend \textbf{une routine} que vous
utiliserez dans tous vos TP et dans votre projet de 2\textsuperscript{e} année :
\begin{center}
\cmd{git pull} $\rightarrow$ travailler $\rightarrow$ \cmd{git status} $\rightarrow$ \cmd{git add}
$\rightarrow$ \cmd{git commit} $\rightarrow$ \cmd{git push}
\end{center}
À la fin du TP, votre travail sera sauvegardé sur GitHub, avec l'historique de chaque modification.

% =====================================================================
\section{Git et GitHub en deux mots}

\begin{itemize}
  \item \textbf{Git} est un logiciel installé sur votre PC. Il enregistre l'historique de vos fichiers : chaque
  sauvegarde s'appelle un \textbf{commit} (une « photo » de vos fichiers, avec un message et votre nom).
  \item \textbf{GitHub} est un site web qui héberge une copie de votre travail. Vous le retrouvez depuis n'importe
  quel PC, et vous pouvez le partager avec votre équipe.
  \item Un \textbf{dépôt} (\emph{repository}, ou \emph{repo}) est un dossier de projet suivi par Git. Il existe en
  deux exemplaires : le dépôt \textbf{distant} sur GitHub et le dépôt \textbf{local} sur votre PC.
\end{itemize}

\begin{center}
\begin{tikzpicture}[node distance=1.55cm]
  \node[zone=blue!8]                  (wd)  {Fichiers\\modifiés};
  \node[zone=orange!12, right=of wd]  (idx) {Prêts à être\\enregistrés};
  \node[zone=yellow!15, right=of idx] (loc) {Historique\\sur le PC};
  \node[zone=green!12, right=of loc]  (gh)  {GitHub\\(en ligne)};
  \draw[->, thick] (wd.north east)  to[bend left=35] node[above, font=\small]{\texttt{git add}}    (idx.north west);
  \draw[->, thick] (idx.north east) to[bend left=35] node[above, font=\small]{\texttt{git commit}} (loc.north west);
  \draw[->, thick] (loc.north east) to[bend left=35] node[above, font=\small]{\texttt{git push}}   (gh.north west);
  \draw[->, thick, dashed] (gh.south) to[bend left=14] node[below=1pt, font=\small]{\texttt{git pull} : récupérer les nouveautés de GitHub} (wd.south);
  \draw[gray, rounded corners, densely dotted] ([xshift=-2mm,yshift=11mm]wd.north west) rectangle ([xshift=2mm,yshift=-3mm]loc.south east);
  \node[font=\footnotesize\itshape, text=gray!60!black, anchor=north west] at ([xshift=-1mm,yshift=10mm]wd.north west) {sur votre PC};
\end{tikzpicture}
\end{center}

\begin{info}[Commit ou push ?]
\cmd{git commit} enregistre une version \emph{sur votre PC}. Tant que vous n'avez pas fait \cmd{git push},
GitHub ne la connaît pas : si le PC tombe en panne ou si vous changez de poste, elle est perdue.
\end{info}

% =====================================================================
\section{Préparer GitHub et son poste}

\begin{exercice}{Créer son compte et son premier dépôt sur GitHub}
\begin{enumerate}
  \item Sur \url{https://github.com}, créez un compte (\gh{Sign up}). Choisissez un nom d'utilisateur sérieux :
  il pourra apparaître sur votre CV. Notez-le. \reponse{1}
  \item Créez un dépôt : bouton \gh{+} en haut à droite, puis \gh{New repository}.
  \begin{itemize}
    \item \emph{Repository name} : \cmd{tp-git} ;
    \item visibilité : \gh{Private} ;
    \item activez l'option qui ajoute un fichier \gh{README} ;
    \item validez avec \gh{Create repository}.
  \end{itemize}
  \item Que contient votre dépôt ? Cliquez sur l'horloge « \gh{1 Commit} » : quel est le message de ce commit ? \reponse{2}
\end{enumerate}
\begin{correction}
\begin{enumerate}
  \item[3.] Un seul fichier, \cmd{README.md}, créé par GitHub avec le commit \emph{Initial commit}.
  Le dépôt a donc déjà un historique d'un commit.
\end{enumerate}
\end{correction}
\end{exercice}

\textbf{Configurer Git sur le poste.} Git inscrit votre nom et votre e-mail dans chaque commit. Ouvrez
\textbf{Git Bash} (Windows) ou un terminal (Linux) et tapez, avec \emph{vos} informations :
\begin{lstlisting}
git --version
git config --global user.name "Prénom Nom"
git config --global user.email "adresse-du-compte-github@exemple.fr"
git config --global core.editor nano
git config --global pull.rebase false
\end{lstlisting}
Les deux dernières lignes règlent le comportement de Git quand il fusionne votre travail avec celui d'un camarade
(partie 6). Vérifiez avec \cmd{git config --global --list}.

\begin{attention}[Poste partagé]
Sur un PC de la salle utilisé par d'autres élèves, vérifiez en début de séance que
\cmd{git config user.name} affiche bien \emph{votre} nom. Sinon, refaites la configuration.
\end{attention}

% =====================================================================
\section{Récupérer son dépôt sur le PC : git clone}

Sur la page GitHub de votre dépôt, cliquez sur le bouton vert \gh{Code}, onglet \gh{HTTPS}, et copiez l'adresse.
Dans le terminal (dans Git Bash, collez avec \textbf{Maj+Inser} ou clic droit $\rightarrow$ \emph{Paste}) :
\begin{lstlisting}
cd ~
git clone https://github.com/VOTRE-COMPTE/tp-git.git
cd tp-git
ls
git status
\end{lstlisting}

\cmd{git clone} ne se fait \textbf{qu'une seule fois} par poste : il crée le dossier \cmd{tp-git}, relié à GitHub.

\begin{info}[Connexion à GitHub]
La première fois, Git demande de vous connecter. Sous Windows, une fenêtre s'ouvre : choisissez
\gh{Sign in with your browser} et autorisez l'accès. Le mot de passe du compte n'est pas accepté directement
dans le terminal. Sous Linux, utilisez un \emph{jeton d'accès personnel} (\emph{personal access token}) créé dans
\gh{Settings} $\rightarrow$ \gh{Developer settings}.
\end{info}

\begin{exercice}{Cloner son dépôt}
\begin{enumerate}
  \item Clonez votre dépôt avec les commandes ci-dessus. Quel(s) fichier(s) voyez-vous avec \cmd{ls} ? \reponse{1}
  \item Recopiez les deux premières lignes affichées par \cmd{git status}. \reponse{2}
\end{enumerate}
\begin{correction}
\begin{enumerate}
  \item \cmd{README.md}. Il existe aussi un dossier caché \cmd{.git} (visible avec \cmd{ls -a}) : c'est là que Git
  range l'historique. On n'y touche jamais.
  \item \cmd{On branch main} puis \cmd{Your branch is up to date with 'origin/main'.} (en français :
  « Sur la branche main. Votre branche est à jour avec 'origin/main'. »). \cmd{origin} désigne GitHub.
  En dernière ligne : \emph{nothing to commit, working tree clean}, rien à enregistrer.
\end{enumerate}
\end{correction}
\end{exercice}

% =====================================================================
\section{La routine : status, add, commit, push}

\begin{tabularx}{\linewidth}{@{}>{\raggedright\arraybackslash}p{5.3cm} X@{}}
\toprule
\textbf{Commande} & \textbf{Rôle} \\
\midrule
\cmd{git status}               & « Où en suis-je ? » : quels fichiers sont nouveaux, modifiés, prêts à être enregistrés. \textbf{À taper tout le temps.} \\
\cmd{git add fichier}          & sélectionne un fichier pour le prochain commit (\cmd{git add .} sélectionne tout le dossier) \\
\cmd{git commit -m "message"}  & enregistre les fichiers sélectionnés, avec un message qui dit ce qui a changé \\
\cmd{git push}                 & envoie les commits vers GitHub \\
\bottomrule
\end{tabularx}

\begin{exercice}{Premier cycle complet (guidé)}
Vous allez documenter le réseau de la salle B204.
\begin{enumerate}
  \item Dans le dossier \cmd{tp-git}, créez le fichier \cmd{plan_adressage.txt} avec \cmd{nano plan_adressage.txt}
  (enregistrer : \textbf{Ctrl+O}, \textbf{Entrée} ; quitter : \textbf{Ctrl+X}) :
\begin{lstlisting}
Plan d'adressage - salle B204
VLAN 10  ADMIN    192.168.10.0/24  passerelle 192.168.10.254
VLAN 20  ELEVES   192.168.20.0/24  passerelle 192.168.20.254
\end{lstlisting}
  \item \cmd{git status} : dans quelle rubrique, et de quelle couleur, apparaît le fichier ? \reponse{1}
  \item \cmd{git add plan_adressage.txt}, puis \cmd{git status}. Qu'est-ce qui a changé ? \reponse{1}
  \item Enregistrez :
\begin{lstlisting}
git commit -m "Ajout du plan d'adressage de la salle B204"
git status
\end{lstlisting}
  Que vous dit maintenant \cmd{git status} ? Quelle commande vous suggère-t-il ? \reponse{2}
  \item Envoyez vers GitHub avec \cmd{git push}, puis \cmd{git status}. Rafraîchissez la page GitHub : que voyez-vous ?
  \reponse{2}
\end{enumerate}
\begin{correction}
\begin{enumerate}
  \item[2.] \emph{Untracked files} (fichiers non suivis), en rouge : Git voit le fichier mais ne le suit pas encore.
  \item[3.] Le fichier passe dans \emph{Changes to be committed}, en vert : il sera dans le prochain commit.
  \item[4.] \cmd{Your branch is ahead of 'origin/main' by 1 commit.} suivi de
  \cmd{(use "git push" to publish your local commits)}. Le commit existe sur le PC, pas encore sur GitHub.
  \item[5.] \cmd{git status} revient à \emph{up to date with 'origin/main'}. Sur GitHub, le fichier
  \cmd{plan_adressage.txt} apparaît avec le message du commit, et le compteur passe à 2 commits.
\end{enumerate}
\end{correction}
\end{exercice}

\begin{info}[Lire git status]
\begin{tabularx}{\linewidth}{@{}>{\raggedright\arraybackslash}p{7cm} X@{}}
\emph{Untracked files} (rouge)            & nouveau fichier, pas encore suivi $\rightarrow$ \cmd{git add} \\
\emph{Changes not staged for commit} (rouge) & fichier modifié, pas encore sélectionné $\rightarrow$ \cmd{git add} \\
\emph{Changes to be committed} (vert)     & prêt à être enregistré $\rightarrow$ \cmd{git commit} \\
\emph{Your branch is ahead of 'origin/main'} & commits pas encore envoyés $\rightarrow$ \cmd{git push} \\
\emph{nothing to commit, working tree clean} & tout est enregistré \\
\end{tabularx}
\end{info}

\begin{exercice}{À vous de jouer (sans aide)}
\begin{enumerate}
  \item Ajoutez au plan d'adressage la ligne suivante :
\begin{lstlisting}
VLAN 30  SERVEURS 192.168.30.0/24  passerelle 192.168.30.254
\end{lstlisting}
  \item Créez un fichier \cmd{inventaire.csv} :
\begin{lstlisting}
nom;type;ip
SW-B204;switch;192.168.10.2
PC-B204-01;poste;192.168.20.11
\end{lstlisting}
  \item Faites \textbf{deux commits séparés} (un par modification), avec des messages clairs, puis envoyez-les sur
  GitHub. Notez ci-dessous toutes les commandes utilisées, dans l'ordre. \reponse{7}
  \item Sur GitHub, cliquez sur le compteur de commits. Combien y en a-t-il ? \reponse{1}
\end{enumerate}
\begin{correction}
\begin{lstlisting}
nano plan_adressage.txt
git status
git add plan_adressage.txt
git commit -m "Ajout du VLAN 30 (serveurs)"
nano inventaire.csv
git add inventaire.csv
git commit -m "Ajout de l'inventaire du matériel"
git push
\end{lstlisting}
Un seul \cmd{git push} suffit pour envoyer les deux commits. GitHub affiche 4 commits : \emph{Initial commit} et les
trois commits de l'élève. On valorise les \cmd{git status} intercalés.
\end{correction}
\end{exercice}

\textbf{Voir l'historique.} Sur le PC, \cmd{git log --oneline} affiche un commit par ligne, du plus récent au plus
ancien. Sur GitHub, le compteur de commits donne la même liste. En cliquant sur un commit, on voit les lignes
ajoutées (en vert) et supprimées (en rouge).

\begin{exercice}{Bon ou mauvais message ?}
Un bon message dit \emph{ce qui a changé}, en une phrase courte. Pour chaque message, indiquez B (bon) ou M (mauvais).

\renewcommand{\arraystretch}{1.5}
\begin{tabularx}{\linewidth}{@{}l X c@{}}
\toprule
& \textbf{Message} & \textbf{B/M} \\
\midrule
1 & « modif » & \ifcorrige M\fi \\
2 & « Ajout du VLAN 30 (serveurs) au plan d'adressage » & \ifcorrige B\fi \\
3 & « azerty » & \ifcorrige M\fi \\
4 & « Correction de l'adresse IP du switch dans l'inventaire » & \ifcorrige B\fi \\
5 & « travail de mardi » & \ifcorrige M\fi \\
\bottomrule
\end{tabularx}
\renewcommand{\arraystretch}{1}
\begin{correction}
Les messages 1, 3 et 5 ne disent rien de la modification : dans trois mois, personne (ni vous) ne saura ce qu'ils
contiennent.
\end{correction}
\end{exercice}

% =====================================================================
\section{Récupérer les nouveautés : git pull}

Quand le dépôt GitHub a changé (modification par un camarade, ou par vous depuis un autre PC), on récupère ces
changements avec \cmd{git pull}.

\begin{exercice}{Modifier sur GitHub, récupérer sur le PC}
\begin{enumerate}
  \item Sur GitHub, ouvrez \cmd{README.md}, cliquez sur le crayon (\gh{Edit this file}) et ajoutez la ligne :
  \emph{Documentation du réseau de la salle B204 -- Prénom Nom}. Validez avec \gh{Commit changes}.
  \item Sur le PC, tapez \cmd{git status}. Git signale-t-il la modification faite sur GitHub ? \reponse{1}
  \item Tapez \cmd{git pull}, puis affichez le README avec \cmd{cat README.md}. \reponse{1}
  \item Pourquoi \cmd{git status} ne voyait-il pas la modification à la question 2 ? \reponse{2}
\end{enumerate}
\begin{correction}
\begin{enumerate}
  \item[2.] Non : il affiche toujours \emph{up to date with 'origin/main'}.
  \item[3.] Git affiche \emph{Fast-forward} et le README contient la nouvelle ligne.
  \item[4.] \cmd{git status} ne contacte pas GitHub : il compare avec le dernier état de GitHub \emph{connu} du PC.
  Seuls \cmd{git pull} (et \cmd{git fetch}) vont chercher les nouveautés. D'où la règle :
  \textbf{toujours commencer par \cmd{git pull}}.
\end{enumerate}
\end{correction}
\end{exercice}

\begin{attention}[Jamais de mot de passe sur GitHub]
Tout ce qui est poussé sur GitHub reste dans l'historique, même si l'on supprime le fichier ensuite.
Ne versionnez jamais un mot de passe, une clé ou une configuration contenant un secret. Si cela arrive,
le mot de passe doit être considéré comme compromis et changé.
\end{attention}

% =====================================================================
\section{Travailler à deux sur le même dépôt}

\begin{exercice}{En binôme (élèves A et B)}
\begin{enumerate}
  \item \textbf{(A)} Sur la page GitHub de votre dépôt : \gh{Settings} $\rightarrow$ \gh{Collaborators}
  $\rightarrow$ \gh{Add people}, et ajoutez le compte de B. \textbf{(B)} Acceptez l'invitation reçue par e-mail
  (ou sur \url{https://github.com/notifications}).
  \item \textbf{(B)} Clonez le dépôt de A dans un nouveau dossier :
\begin{lstlisting}
cd ~
git clone https://github.com/COMPTE-DE-A/tp-git.git tp-git-A
cd tp-git-A
\end{lstlisting}
  \item \textbf{(B)} Ajoutez la ligne \cmd{SW-B204-2;switch;192.168.10.3} à la fin de \cmd{inventaire.csv}, puis
  faites la routine jusqu'au \cmd{git push}.
  \item \textbf{(A)} Faites \cmd{git pull} et vérifiez que la ligne de B est bien dans votre inventaire. \reponse{1}
  \item \textbf{En même temps} : A ajoute une ligne au README, B ajoute un VLAN au plan d'adressage. Chacun fait
  \cmd{add} et \cmd{commit}. B fait \cmd{git push} en premier, \emph{puis} A. Que se passe-t-il pour A ? \reponse{2}
  \item \textbf{(A)} Faites \cmd{git pull}. Si \cmd{nano} s'ouvre avec un message de fusion, enregistrez
  (\textbf{Ctrl+O}, \textbf{Entrée}) et quittez (\textbf{Ctrl+X}). Puis \cmd{git push}. Tout le monde fait ensuite
  \cmd{git pull} : les deux dépôts sont-ils identiques ? \reponse{1}
\end{enumerate}
\begin{correction}
\begin{enumerate}
  \item[4.] La ligne de B apparaît chez A.
  \item[5.] Le \cmd{git push} de A est refusé : \cmd{! [rejected] main -> main (fetch first)}. GitHub contient un
  commit (celui de B) que A n'a pas encore. Rien n'est perdu.
  \item[6.] \cmd{git pull} récupère le commit de B et le fusionne avec celui de A (fichiers différents : fusion
  automatique). Après le \cmd{git push} de A et le \cmd{git pull} de B, les deux dépôts sont identiques.
\end{enumerate}
Si les deux élèves modifient \emph{la même ligne}, Git signale un \emph{CONFLICT} : hors programme de ce TP,
l'enseignant aide à le résoudre (garder la bonne ligne, supprimer les marqueurs \cmd{<<<<<<<}, \cmd{=======},
\cmd{>>>>>>>}, puis \cmd{add}, \cmd{commit}, \cmd{push}).
\end{correction}
\end{exercice}

\begin{info}[« push rejected » : pas de panique]
Si \cmd{git push} est refusé, c'est que GitHub a des nouveautés que vous n'avez pas. Faites \cmd{git pull}, puis
\cmd{git push}. Si Git parle de \emph{CONFLICT}, appelez l'enseignant.
\end{info}

% =====================================================================
\section{Défi final (en autonomie)}

Sans regarder les pages précédentes, en vous aidant seulement de l'aide-mémoire :
\begin{enumerate}
  \item Créez sur GitHub un nouveau dépôt privé \cmd{mes-tp-ciel}, avec un README, et clonez-le.
  \item Ajoutez un fichier \cmd{compte_rendu.md} contenant le titre de ce TP et trois choses que vous avez apprises.
  Commit et push.
  \item Complétez le README avec votre nom et votre classe. Commit et push.
  \item Modifiez \cmd{compte_rendu.md} directement sur GitHub, puis récupérez la modification sur le PC.
  \item Montrez à l'enseignant la liste des commits sur GitHub et un \cmd{git status} « propre ».
\end{enumerate}
Ce dépôt vous servira pour ranger vos prochains TP.

\begin{correction}
Attendus : au moins 3 commits (\emph{Initial commit}, \cmd{compte_rendu.md}, README) plus le commit fait sur GitHub,
visible sur le PC après \cmd{git pull}. \cmd{git status} doit afficher \emph{up to date with 'origin/main'} et
\emph{nothing to commit, working tree clean}. Vérifier aussi la qualité des messages.
\end{correction}

\subsection*{Auto-évaluation}
\begin{tabularx}{\linewidth}{@{}X c c c@{}}
\toprule
\textbf{Je sais\ldots} & \textbf{Pas encore} & \textbf{En partie} & \textbf{Oui} \\
\midrule
expliquer la différence entre Git et GitHub, entre commit et push & $\square$ & $\square$ & $\square$ \\
créer un dépôt sur GitHub et le cloner sur un PC & $\square$ & $\square$ & $\square$ \\
lire \cmd{git status} et savoir quelle commande taper ensuite & $\square$ & $\square$ & $\square$ \\
enregistrer mon travail : \cmd{add}, \cmd{commit} avec un message clair & $\square$ & $\square$ & $\square$ \\
envoyer et récupérer mon travail : \cmd{push}, \cmd{pull} & $\square$ & $\square$ & $\square$ \\
réagir à un \emph{push rejected} & $\square$ & $\square$ & $\square$ \\
\bottomrule
\end{tabularx}

% =====================================================================
\vspace{5mm}
\needspace{20\baselineskip}
\section*{Aide-mémoire}
\addcontentsline{toc}{section}{Aide-mémoire}

\begin{tcolorbox}[colback=ciel!5, colframe=ciel, title=La routine de chaque séance, fonttitle=\bfseries]
\begin{tabularx}{\linewidth}{@{}l >{\raggedright\arraybackslash}p{5.5cm} X@{}}
\textbf{1.} & \cmd{git pull}                 & récupérer les nouveautés (début de séance) \\
\textbf{2.} & \emph{travailler}              & créer, modifier des fichiers \\
\textbf{3.} & \cmd{git status}               & voir ce qui a changé \\
\textbf{4.} & \cmd{git add fichier}          & sélectionner (ou \cmd{git add .} pour tout) \\
\textbf{5.} & \cmd{git commit -m "message"}  & enregistrer sur le PC \\
\textbf{6.} & \cmd{git push}                 & envoyer sur GitHub (au plus tard en fin de séance) \\
\end{tabularx}
\end{tcolorbox}

\begin{tabularx}{\linewidth}{@{}>{\raggedright\arraybackslash}p{7.3cm} X@{}}
\toprule
\multicolumn{2}{@{}l}{\textbf{\color{ciel}Une seule fois}} \\
\cmd{git config --global user.name "Nom"} & votre nom (dans chaque commit) \\
\cmd{git config --global user.email "mail"} & votre e-mail (celui du compte GitHub) \\
\cmd{git clone <url>} & copier un dépôt GitHub sur le PC \\
\midrule
\multicolumn{2}{@{}l}{\textbf{\color{ciel}Pour s'y retrouver}} \\
\cmd{git status} & où en suis-je ? \\
\cmd{git log --oneline} & liste des commits \\
\bottomrule
\end{tabularx}

\subsection*{Pour aller plus loin}
\begin{itemize}
  \item S. Chacon et B. Straub, \emph{Pro Git}, 2\textsuperscript{e} éd., Apress, 2014. Libre, traduit en français :
  \url{https://git-scm.com/book/fr/v2} (chapitres 1 et 2).
  \item Documentation GitHub en français : \url{https://docs.github.com/fr/get-started}.
\end{itemize}

\end{document}
"
%  Compiler deux fois (numéro de dernière page dans le pied de page).
% =====================================================================
\documentclass[11pt,a4paper]{article}

\usepackage[utf8]{inputenc}
\usepackage[T1]{fontenc}
\usepackage[french]{babel}
\usepackage{lmodern}
\usepackage{microtype}
% Pas de ligatures -- et << >> dans les commandes (police à chasse fixe)
\DisableLigatures[-,<,>]{encoding=T1, family=tt*}
\usepackage{needspace}
\usepackage[margin=2.2cm,headheight=14pt]{geometry}
\usepackage{xcolor}
\usepackage{upquote}
\usepackage{listings}
\usepackage[most]{tcolorbox}
\usepackage{tikz}
\usetikzlibrary{positioning,arrows.meta,babel}
\usepackage{enumitem}
\usepackage{tabularx}
\usepackage{booktabs}
\usepackage{amssymb}
\usepackage{fancyhdr}
\usepackage{lastpage}
\usepackage{comment}
\usepackage[hidelinks]{hyperref}

% ---------------------------------------------------------------------
% Version élève / version corrigée
% ---------------------------------------------------------------------
\newif\ifcorrige
\ifdefined\CORRIGE\corrigetrue\fi

\definecolor{ciel}{RGB}{0,84,147}

\ifcorrige
  \newtcolorbox{correction}{breakable, colback=green!4, colframe=green!45!black,
    fonttitle=\bfseries\small, title=Correction, before skip=2mm, after skip=3mm}
  \newcommand{\reponse}[1]{}
\else
  \excludecomment{correction}
  \newcommand{\reponse}[1]{\par\smallskip{\color{gray!70}%
    \foreach \n in {1,...,#1}{\noindent\mbox{}\dotfill\par\vspace{2.5mm}}}}
\fi

% ---------------------------------------------------------------------
% Mise en forme
% ---------------------------------------------------------------------
\setlength{\parindent}{0pt}
\setlength{\parskip}{4pt}
\setlist{itemsep=2pt, topsep=3pt}

\pagestyle{fancy}
\fancyhf{}
\lhead{\small BTS CIEL 1 -- Gestion de projet}
\rhead{\small TP Git \& GitHub\ifcorrige{} -- \textcolor{red!70!black}{corrigé}\fi}
\cfoot{\small \thepage/\pageref{LastPage}}

% Commande en ligne : \detokenize permet d'écrire _ ~ : sans échappement
% (pas d'accents ni de # dans l'argument).
\newcommand{\cmd}[1]{\texttt{\textcolor{ciel!70!black}{\detokenize{#1}}}}
% Élément de l'interface web de GitHub
\newcommand{\gh}[1]{\textsf{\textbf{#1}}}

\lstdefinestyle{bash}{
  basicstyle=\ttfamily\small, backgroundcolor=\color{gray!8},
  frame=single, rulecolor=\color{gray!40}, framesep=4pt,
  xleftmargin=4pt, xrightmargin=4pt,
  columns=fullflexible, keepspaces=true, breaklines=true,
  showstringspaces=false, upquote=true, aboveskip=6pt, belowskip=6pt,
  literate={é}{{\'e}}1 {è}{{\`e}}1 {ê}{{\^e}}1 {à}{{\`a}}1 {â}{{\^a}}1
           {ù}{{\`u}}1 {û}{{\^u}}1 {ô}{{\^o}}1 {î}{{\^i}}1 {ç}{{\c{c}}}1
           {É}{{\'E}}1 {È}{{\`E}}1 {À}{{\`A}}1
}
\lstdefinestyle{sortie}{style=bash, backgroundcolor=\color{white},
  frame=leftline, framerule=2pt, rulecolor=\color{ciel!50}, basicstyle=\ttfamily\small\color{black!75}}
\lstset{style=bash}

\newtcolorbox{info}[1][À retenir]{breakable, colback=ciel!5, colframe=ciel!80!black,
  fonttitle=\bfseries, title=#1}
\newtcolorbox{attention}[1][Attention]{breakable, colback=red!4, colframe=red!60!black,
  fonttitle=\bfseries, title=#1}
\newtcolorbox[auto counter]{exercice}[1]{breakable, colback=white, colframe=orange!80!black,
  fonttitle=\bfseries, title={Exercice~\thetcbcounter~--~#1}}

\tikzset{
  zone/.style={draw, thick, rounded corners, minimum width=2.9cm, minimum height=1.4cm,
               align=center, fill=#1, font=\small},
  >={Stealth[length=2.5mm]}
}

\begin{document}

% =====================================================================
\begin{tcolorbox}[colback=ciel!6, colframe=ciel, arc=2mm]
  {\LARGE\bfseries TP -- Premiers pas avec Git et GitHub}\\[1mm]
  {\large Sauvegarder et partager son travail, la bonne routine}\\[2mm]
  BTS CIEL -- 1\textsuperscript{re} année \hfill Gestion de projet
  \ifcorrige\\[2mm]{\bfseries\color{red!70!black} VERSION CORRIGÉE -- DOCUMENT ENSEIGNANT}\fi
  \tcblower
  \begin{tabularx}{\linewidth}{@{}>{\bfseries}l X@{}}
    Durée      & 2 h à 3 h \\
    Prérequis  & ouvrir un terminal, utiliser \cmd{cd} et \cmd{ls}, éditer un fichier avec \cmd{nano} ou VS Code \\
    Matériel   & poste avec Git (Git Bash sous Windows), navigateur web, adresse e-mail pour créer un compte GitHub \\
  \end{tabularx}
\end{tcolorbox}

\ifcorrige\else
\vspace{2mm}
\textbf{Nom :} \makebox[6cm]{\dotfill} \hfill \textbf{Prénom :} \makebox[6cm]{\dotfill}
\vspace{2mm}
\fi

\textbf{Objectif.} Ce TP ne fait pas de vous des experts de Git. Il vous apprend \textbf{une routine} que vous
utiliserez dans tous vos TP et dans votre projet de 2\textsuperscript{e} année :
\begin{center}
\cmd{git pull} $\rightarrow$ travailler $\rightarrow$ \cmd{git status} $\rightarrow$ \cmd{git add}
$\rightarrow$ \cmd{git commit} $\rightarrow$ \cmd{git push}
\end{center}
À la fin du TP, votre travail sera sauvegardé sur GitHub, avec l'historique de chaque modification.

% =====================================================================
\section{Git et GitHub en deux mots}

\begin{itemize}
  \item \textbf{Git} est un logiciel installé sur votre PC. Il enregistre l'historique de vos fichiers : chaque
  sauvegarde s'appelle un \textbf{commit} (une « photo » de vos fichiers, avec un message et votre nom).
  \item \textbf{GitHub} est un site web qui héberge une copie de votre travail. Vous le retrouvez depuis n'importe
  quel PC, et vous pouvez le partager avec votre équipe.
  \item Un \textbf{dépôt} (\emph{repository}, ou \emph{repo}) est un dossier de projet suivi par Git. Il existe en
  deux exemplaires : le dépôt \textbf{distant} sur GitHub et le dépôt \textbf{local} sur votre PC.
\end{itemize}

\begin{center}
\begin{tikzpicture}[node distance=1.55cm]
  \node[zone=blue!8]                  (wd)  {Fichiers\\modifiés};
  \node[zone=orange!12, right=of wd]  (idx) {Prêts à être\\enregistrés};
  \node[zone=yellow!15, right=of idx] (loc) {Historique\\sur le PC};
  \node[zone=green!12, right=of loc]  (gh)  {GitHub\\(en ligne)};
  \draw[->, thick] (wd.north east)  to[bend left=35] node[above, font=\small]{\texttt{git add}}    (idx.north west);
  \draw[->, thick] (idx.north east) to[bend left=35] node[above, font=\small]{\texttt{git commit}} (loc.north west);
  \draw[->, thick] (loc.north east) to[bend left=35] node[above, font=\small]{\texttt{git push}}   (gh.north west);
  \draw[->, thick, dashed] (gh.south) to[bend left=14] node[below=1pt, font=\small]{\texttt{git pull} : récupérer les nouveautés de GitHub} (wd.south);
  \draw[gray, rounded corners, densely dotted] ([xshift=-2mm,yshift=11mm]wd.north west) rectangle ([xshift=2mm,yshift=-3mm]loc.south east);
  \node[font=\footnotesize\itshape, text=gray!60!black, anchor=north west] at ([xshift=-1mm,yshift=10mm]wd.north west) {sur votre PC};
\end{tikzpicture}
\end{center}

\begin{info}[Commit ou push ?]
\cmd{git commit} enregistre une version \emph{sur votre PC}. Tant que vous n'avez pas fait \cmd{git push},
GitHub ne la connaît pas : si le PC tombe en panne ou si vous changez de poste, elle est perdue.
\end{info}

% =====================================================================
\section{Préparer GitHub et son poste}

\begin{exercice}{Créer son compte et son premier dépôt sur GitHub}
\begin{enumerate}
  \item Sur \url{https://github.com}, créez un compte (\gh{Sign up}). Choisissez un nom d'utilisateur sérieux :
  il pourra apparaître sur votre CV. Notez-le. \reponse{1}
  \item Créez un dépôt : bouton \gh{+} en haut à droite, puis \gh{New repository}.
  \begin{itemize}
    \item \emph{Repository name} : \cmd{tp-git} ;
    \item visibilité : \gh{Private} ;
    \item activez l'option qui ajoute un fichier \gh{README} ;
    \item validez avec \gh{Create repository}.
  \end{itemize}
  \item Que contient votre dépôt ? Cliquez sur l'horloge « \gh{1 Commit} » : quel est le message de ce commit ? \reponse{2}
\end{enumerate}
\begin{correction}
\begin{enumerate}
  \item[3.] Un seul fichier, \cmd{README.md}, créé par GitHub avec le commit \emph{Initial commit}.
  Le dépôt a donc déjà un historique d'un commit.
\end{enumerate}
\end{correction}
\end{exercice}

\textbf{Configurer Git sur le poste.} Git inscrit votre nom et votre e-mail dans chaque commit. Ouvrez
\textbf{Git Bash} (Windows) ou un terminal (Linux) et tapez, avec \emph{vos} informations :
\begin{lstlisting}
git --version
git config --global user.name "Prénom Nom"
git config --global user.email "adresse-du-compte-github@exemple.fr"
git config --global core.editor nano
git config --global pull.rebase false
\end{lstlisting}
Les deux dernières lignes règlent le comportement de Git quand il fusionne votre travail avec celui d'un camarade
(partie 6). Vérifiez avec \cmd{git config --global --list}.

\begin{attention}[Poste partagé]
Sur un PC de la salle utilisé par d'autres élèves, vérifiez en début de séance que
\cmd{git config user.name} affiche bien \emph{votre} nom. Sinon, refaites la configuration.
\end{attention}

% =====================================================================
\section{Récupérer son dépôt sur le PC : git clone}

Sur la page GitHub de votre dépôt, cliquez sur le bouton vert \gh{Code}, onglet \gh{HTTPS}, et copiez l'adresse.
Dans le terminal (dans Git Bash, collez avec \textbf{Maj+Inser} ou clic droit $\rightarrow$ \emph{Paste}) :
\begin{lstlisting}
cd ~
git clone https://github.com/VOTRE-COMPTE/tp-git.git
cd tp-git
ls
git status
\end{lstlisting}

\cmd{git clone} ne se fait \textbf{qu'une seule fois} par poste : il crée le dossier \cmd{tp-git}, relié à GitHub.

\begin{info}[Connexion à GitHub]
La première fois, Git demande de vous connecter. Sous Windows, une fenêtre s'ouvre : choisissez
\gh{Sign in with your browser} et autorisez l'accès. Le mot de passe du compte n'est pas accepté directement
dans le terminal. Sous Linux, utilisez un \emph{jeton d'accès personnel} (\emph{personal access token}) créé dans
\gh{Settings} $\rightarrow$ \gh{Developer settings}.
\end{info}

\begin{exercice}{Cloner son dépôt}
\begin{enumerate}
  \item Clonez votre dépôt avec les commandes ci-dessus. Quel(s) fichier(s) voyez-vous avec \cmd{ls} ? \reponse{1}
  \item Recopiez les deux premières lignes affichées par \cmd{git status}. \reponse{2}
\end{enumerate}
\begin{correction}
\begin{enumerate}
  \item \cmd{README.md}. Il existe aussi un dossier caché \cmd{.git} (visible avec \cmd{ls -a}) : c'est là que Git
  range l'historique. On n'y touche jamais.
  \item \cmd{On branch main} puis \cmd{Your branch is up to date with 'origin/main'.} (en français :
  « Sur la branche main. Votre branche est à jour avec 'origin/main'. »). \cmd{origin} désigne GitHub.
  En dernière ligne : \emph{nothing to commit, working tree clean}, rien à enregistrer.
\end{enumerate}
\end{correction}
\end{exercice}

% =====================================================================
\section{La routine : status, add, commit, push}

\begin{tabularx}{\linewidth}{@{}>{\raggedright\arraybackslash}p{5.3cm} X@{}}
\toprule
\textbf{Commande} & \textbf{Rôle} \\
\midrule
\cmd{git status}               & « Où en suis-je ? » : quels fichiers sont nouveaux, modifiés, prêts à être enregistrés. \textbf{À taper tout le temps.} \\
\cmd{git add fichier}          & sélectionne un fichier pour le prochain commit (\cmd{git add .} sélectionne tout le dossier) \\
\cmd{git commit -m "message"}  & enregistre les fichiers sélectionnés, avec un message qui dit ce qui a changé \\
\cmd{git push}                 & envoie les commits vers GitHub \\
\bottomrule
\end{tabularx}

\begin{exercice}{Premier cycle complet (guidé)}
Vous allez documenter le réseau de la salle B204.
\begin{enumerate}
  \item Dans le dossier \cmd{tp-git}, créez le fichier \cmd{plan_adressage.txt} avec \cmd{nano plan_adressage.txt}
  (enregistrer : \textbf{Ctrl+O}, \textbf{Entrée} ; quitter : \textbf{Ctrl+X}) :
\begin{lstlisting}
Plan d'adressage - salle B204
VLAN 10  ADMIN    192.168.10.0/24  passerelle 192.168.10.254
VLAN 20  ELEVES   192.168.20.0/24  passerelle 192.168.20.254
\end{lstlisting}
  \item \cmd{git status} : dans quelle rubrique, et de quelle couleur, apparaît le fichier ? \reponse{1}
  \item \cmd{git add plan_adressage.txt}, puis \cmd{git status}. Qu'est-ce qui a changé ? \reponse{1}
  \item Enregistrez :
\begin{lstlisting}
git commit -m "Ajout du plan d'adressage de la salle B204"
git status
\end{lstlisting}
  Que vous dit maintenant \cmd{git status} ? Quelle commande vous suggère-t-il ? \reponse{2}
  \item Envoyez vers GitHub avec \cmd{git push}, puis \cmd{git status}. Rafraîchissez la page GitHub : que voyez-vous ?
  \reponse{2}
\end{enumerate}
\begin{correction}
\begin{enumerate}
  \item[2.] \emph{Untracked files} (fichiers non suivis), en rouge : Git voit le fichier mais ne le suit pas encore.
  \item[3.] Le fichier passe dans \emph{Changes to be committed}, en vert : il sera dans le prochain commit.
  \item[4.] \cmd{Your branch is ahead of 'origin/main' by 1 commit.} suivi de
  \cmd{(use "git push" to publish your local commits)}. Le commit existe sur le PC, pas encore sur GitHub.
  \item[5.] \cmd{git status} revient à \emph{up to date with 'origin/main'}. Sur GitHub, le fichier
  \cmd{plan_adressage.txt} apparaît avec le message du commit, et le compteur passe à 2 commits.
\end{enumerate}
\end{correction}
\end{exercice}

\begin{info}[Lire git status]
\begin{tabularx}{\linewidth}{@{}>{\raggedright\arraybackslash}p{7cm} X@{}}
\emph{Untracked files} (rouge)            & nouveau fichier, pas encore suivi $\rightarrow$ \cmd{git add} \\
\emph{Changes not staged for commit} (rouge) & fichier modifié, pas encore sélectionné $\rightarrow$ \cmd{git add} \\
\emph{Changes to be committed} (vert)     & prêt à être enregistré $\rightarrow$ \cmd{git commit} \\
\emph{Your branch is ahead of 'origin/main'} & commits pas encore envoyés $\rightarrow$ \cmd{git push} \\
\emph{nothing to commit, working tree clean} & tout est enregistré \\
\end{tabularx}
\end{info}

\begin{exercice}{À vous de jouer (sans aide)}
\begin{enumerate}
  \item Ajoutez au plan d'adressage la ligne suivante :
\begin{lstlisting}
VLAN 30  SERVEURS 192.168.30.0/24  passerelle 192.168.30.254
\end{lstlisting}
  \item Créez un fichier \cmd{inventaire.csv} :
\begin{lstlisting}
nom;type;ip
SW-B204;switch;192.168.10.2
PC-B204-01;poste;192.168.20.11
\end{lstlisting}
  \item Faites \textbf{deux commits séparés} (un par modification), avec des messages clairs, puis envoyez-les sur
  GitHub. Notez ci-dessous toutes les commandes utilisées, dans l'ordre. \reponse{7}
  \item Sur GitHub, cliquez sur le compteur de commits. Combien y en a-t-il ? \reponse{1}
\end{enumerate}
\begin{correction}
\begin{lstlisting}
nano plan_adressage.txt
git status
git add plan_adressage.txt
git commit -m "Ajout du VLAN 30 (serveurs)"
nano inventaire.csv
git add inventaire.csv
git commit -m "Ajout de l'inventaire du matériel"
git push
\end{lstlisting}
Un seul \cmd{git push} suffit pour envoyer les deux commits. GitHub affiche 4 commits : \emph{Initial commit} et les
trois commits de l'élève. On valorise les \cmd{git status} intercalés.
\end{correction}
\end{exercice}

\textbf{Voir l'historique.} Sur le PC, \cmd{git log --oneline} affiche un commit par ligne, du plus récent au plus
ancien. Sur GitHub, le compteur de commits donne la même liste. En cliquant sur un commit, on voit les lignes
ajoutées (en vert) et supprimées (en rouge).

\begin{exercice}{Bon ou mauvais message ?}
Un bon message dit \emph{ce qui a changé}, en une phrase courte. Pour chaque message, indiquez B (bon) ou M (mauvais).

\renewcommand{\arraystretch}{1.5}
\begin{tabularx}{\linewidth}{@{}l X c@{}}
\toprule
& \textbf{Message} & \textbf{B/M} \\
\midrule
1 & « modif » & \ifcorrige M\fi \\
2 & « Ajout du VLAN 30 (serveurs) au plan d'adressage » & \ifcorrige B\fi \\
3 & « azerty » & \ifcorrige M\fi \\
4 & « Correction de l'adresse IP du switch dans l'inventaire » & \ifcorrige B\fi \\
5 & « travail de mardi » & \ifcorrige M\fi \\
\bottomrule
\end{tabularx}
\renewcommand{\arraystretch}{1}
\begin{correction}
Les messages 1, 3 et 5 ne disent rien de la modification : dans trois mois, personne (ni vous) ne saura ce qu'ils
contiennent.
\end{correction}
\end{exercice}

% =====================================================================
\section{Récupérer les nouveautés : git pull}

Quand le dépôt GitHub a changé (modification par un camarade, ou par vous depuis un autre PC), on récupère ces
changements avec \cmd{git pull}.

\begin{exercice}{Modifier sur GitHub, récupérer sur le PC}
\begin{enumerate}
  \item Sur GitHub, ouvrez \cmd{README.md}, cliquez sur le crayon (\gh{Edit this file}) et ajoutez la ligne :
  \emph{Documentation du réseau de la salle B204 -- Prénom Nom}. Validez avec \gh{Commit changes}.
  \item Sur le PC, tapez \cmd{git status}. Git signale-t-il la modification faite sur GitHub ? \reponse{1}
  \item Tapez \cmd{git pull}, puis affichez le README avec \cmd{cat README.md}. \reponse{1}
  \item Pourquoi \cmd{git status} ne voyait-il pas la modification à la question 2 ? \reponse{2}
\end{enumerate}
\begin{correction}
\begin{enumerate}
  \item[2.] Non : il affiche toujours \emph{up to date with 'origin/main'}.
  \item[3.] Git affiche \emph{Fast-forward} et le README contient la nouvelle ligne.
  \item[4.] \cmd{git status} ne contacte pas GitHub : il compare avec le dernier état de GitHub \emph{connu} du PC.
  Seuls \cmd{git pull} (et \cmd{git fetch}) vont chercher les nouveautés. D'où la règle :
  \textbf{toujours commencer par \cmd{git pull}}.
\end{enumerate}
\end{correction}
\end{exercice}

\begin{attention}[Jamais de mot de passe sur GitHub]
Tout ce qui est poussé sur GitHub reste dans l'historique, même si l'on supprime le fichier ensuite.
Ne versionnez jamais un mot de passe, une clé ou une configuration contenant un secret. Si cela arrive,
le mot de passe doit être considéré comme compromis et changé.
\end{attention}

% =====================================================================
\section{Travailler à deux sur le même dépôt}

\begin{exercice}{En binôme (élèves A et B)}
\begin{enumerate}
  \item \textbf{(A)} Sur la page GitHub de votre dépôt : \gh{Settings} $\rightarrow$ \gh{Collaborators}
  $\rightarrow$ \gh{Add people}, et ajoutez le compte de B. \textbf{(B)} Acceptez l'invitation reçue par e-mail
  (ou sur \url{https://github.com/notifications}).
  \item \textbf{(B)} Clonez le dépôt de A dans un nouveau dossier :
\begin{lstlisting}
cd ~
git clone https://github.com/COMPTE-DE-A/tp-git.git tp-git-A
cd tp-git-A
\end{lstlisting}
  \item \textbf{(B)} Ajoutez la ligne \cmd{SW-B204-2;switch;192.168.10.3} à la fin de \cmd{inventaire.csv}, puis
  faites la routine jusqu'au \cmd{git push}.
  \item \textbf{(A)} Faites \cmd{git pull} et vérifiez que la ligne de B est bien dans votre inventaire. \reponse{1}
  \item \textbf{En même temps} : A ajoute une ligne au README, B ajoute un VLAN au plan d'adressage. Chacun fait
  \cmd{add} et \cmd{commit}. B fait \cmd{git push} en premier, \emph{puis} A. Que se passe-t-il pour A ? \reponse{2}
  \item \textbf{(A)} Faites \cmd{git pull}. Si \cmd{nano} s'ouvre avec un message de fusion, enregistrez
  (\textbf{Ctrl+O}, \textbf{Entrée}) et quittez (\textbf{Ctrl+X}). Puis \cmd{git push}. Tout le monde fait ensuite
  \cmd{git pull} : les deux dépôts sont-ils identiques ? \reponse{1}
\end{enumerate}
\begin{correction}
\begin{enumerate}
  \item[4.] La ligne de B apparaît chez A.
  \item[5.] Le \cmd{git push} de A est refusé : \cmd{! [rejected] main -> main (fetch first)}. GitHub contient un
  commit (celui de B) que A n'a pas encore. Rien n'est perdu.
  \item[6.] \cmd{git pull} récupère le commit de B et le fusionne avec celui de A (fichiers différents : fusion
  automatique). Après le \cmd{git push} de A et le \cmd{git pull} de B, les deux dépôts sont identiques.
\end{enumerate}
Si les deux élèves modifient \emph{la même ligne}, Git signale un \emph{CONFLICT} : hors programme de ce TP,
l'enseignant aide à le résoudre (garder la bonne ligne, supprimer les marqueurs \cmd{<<<<<<<}, \cmd{=======},
\cmd{>>>>>>>}, puis \cmd{add}, \cmd{commit}, \cmd{push}).
\end{correction}
\end{exercice}

\begin{info}[« push rejected » : pas de panique]
Si \cmd{git push} est refusé, c'est que GitHub a des nouveautés que vous n'avez pas. Faites \cmd{git pull}, puis
\cmd{git push}. Si Git parle de \emph{CONFLICT}, appelez l'enseignant.
\end{info}

% =====================================================================
\section{Défi final (en autonomie)}

Sans regarder les pages précédentes, en vous aidant seulement de l'aide-mémoire :
\begin{enumerate}
  \item Créez sur GitHub un nouveau dépôt privé \cmd{mes-tp-ciel}, avec un README, et clonez-le.
  \item Ajoutez un fichier \cmd{compte_rendu.md} contenant le titre de ce TP et trois choses que vous avez apprises.
  Commit et push.
  \item Complétez le README avec votre nom et votre classe. Commit et push.
  \item Modifiez \cmd{compte_rendu.md} directement sur GitHub, puis récupérez la modification sur le PC.
  \item Montrez à l'enseignant la liste des commits sur GitHub et un \cmd{git status} « propre ».
\end{enumerate}
Ce dépôt vous servira pour ranger vos prochains TP.

\begin{correction}
Attendus : au moins 3 commits (\emph{Initial commit}, \cmd{compte_rendu.md}, README) plus le commit fait sur GitHub,
visible sur le PC après \cmd{git pull}. \cmd{git status} doit afficher \emph{up to date with 'origin/main'} et
\emph{nothing to commit, working tree clean}. Vérifier aussi la qualité des messages.
\end{correction}

\subsection*{Auto-évaluation}
\begin{tabularx}{\linewidth}{@{}X c c c@{}}
\toprule
\textbf{Je sais\ldots} & \textbf{Pas encore} & \textbf{En partie} & \textbf{Oui} \\
\midrule
expliquer la différence entre Git et GitHub, entre commit et push & $\square$ & $\square$ & $\square$ \\
créer un dépôt sur GitHub et le cloner sur un PC & $\square$ & $\square$ & $\square$ \\
lire \cmd{git status} et savoir quelle commande taper ensuite & $\square$ & $\square$ & $\square$ \\
enregistrer mon travail : \cmd{add}, \cmd{commit} avec un message clair & $\square$ & $\square$ & $\square$ \\
envoyer et récupérer mon travail : \cmd{push}, \cmd{pull} & $\square$ & $\square$ & $\square$ \\
réagir à un \emph{push rejected} & $\square$ & $\square$ & $\square$ \\
\bottomrule
\end{tabularx}

% =====================================================================
\vspace{5mm}
\needspace{20\baselineskip}
\section*{Aide-mémoire}
\addcontentsline{toc}{section}{Aide-mémoire}

\begin{tcolorbox}[colback=ciel!5, colframe=ciel, title=La routine de chaque séance, fonttitle=\bfseries]
\begin{tabularx}{\linewidth}{@{}l >{\raggedright\arraybackslash}p{5.5cm} X@{}}
\textbf{1.} & \cmd{git pull}                 & récupérer les nouveautés (début de séance) \\
\textbf{2.} & \emph{travailler}              & créer, modifier des fichiers \\
\textbf{3.} & \cmd{git status}               & voir ce qui a changé \\
\textbf{4.} & \cmd{git add fichier}          & sélectionner (ou \cmd{git add .} pour tout) \\
\textbf{5.} & \cmd{git commit -m "message"}  & enregistrer sur le PC \\
\textbf{6.} & \cmd{git push}                 & envoyer sur GitHub (au plus tard en fin de séance) \\
\end{tabularx}
\end{tcolorbox}

\begin{tabularx}{\linewidth}{@{}>{\raggedright\arraybackslash}p{7.3cm} X@{}}
\toprule
\multicolumn{2}{@{}l}{\textbf{\color{ciel}Une seule fois}} \\
\cmd{git config --global user.name "Nom"} & votre nom (dans chaque commit) \\
\cmd{git config --global user.email "mail"} & votre e-mail (celui du compte GitHub) \\
\cmd{git clone <url>} & copier un dépôt GitHub sur le PC \\
\midrule
\multicolumn{2}{@{}l}{\textbf{\color{ciel}Pour s'y retrouver}} \\
\cmd{git status} & où en suis-je ? \\
\cmd{git log --oneline} & liste des commits \\
\bottomrule
\end{tabularx}

\subsection*{Pour aller plus loin}
\begin{itemize}
  \item S. Chacon et B. Straub, \emph{Pro Git}, 2\textsuperscript{e} éd., Apress, 2014. Libre, traduit en français :
  \url{https://git-scm.com/book/fr/v2} (chapitres 1 et 2).
  \item Documentation GitHub en français : \url{https://docs.github.com/fr/get-started}.
\end{itemize}

\end{document}
"
%  Compiler deux fois (numéro de dernière page dans le pied de page).
% =====================================================================
\documentclass[11pt,a4paper]{article}

\usepackage[utf8]{inputenc}
\usepackage[T1]{fontenc}
\usepackage[french]{babel}
\usepackage{lmodern}
\usepackage{microtype}
% Pas de ligatures -- et << >> dans les commandes (police à chasse fixe)
\DisableLigatures[-,<,>]{encoding=T1, family=tt*}
\usepackage{needspace}
\usepackage[margin=2.2cm,headheight=14pt]{geometry}
\usepackage{xcolor}
\usepackage{upquote}
\usepackage{listings}
\usepackage[most]{tcolorbox}
\usepackage{tikz}
\usetikzlibrary{positioning,arrows.meta,babel}
\usepackage{enumitem}
\usepackage{tabularx}
\usepackage{booktabs}
\usepackage{amssymb}
\usepackage{fancyhdr}
\usepackage{lastpage}
\usepackage{comment}
\usepackage[hidelinks]{hyperref}

% ---------------------------------------------------------------------
% Version élève / version corrigée
% ---------------------------------------------------------------------
\newif\ifcorrige
\ifdefined\CORRIGE\corrigetrue\fi

\definecolor{ciel}{RGB}{0,84,147}

\ifcorrige
  \newtcolorbox{correction}{breakable, colback=green!4, colframe=green!45!black,
    fonttitle=\bfseries\small, title=Correction, before skip=2mm, after skip=3mm}
  \newcommand{\reponse}[1]{}
\else
  \excludecomment{correction}
  \newcommand{\reponse}[1]{\par\smallskip{\color{gray!70}%
    \foreach \n in {1,...,#1}{\noindent\mbox{}\dotfill\par\vspace{2.5mm}}}}
\fi

% ---------------------------------------------------------------------
% Mise en forme
% ---------------------------------------------------------------------
\setlength{\parindent}{0pt}
\setlength{\parskip}{4pt}
\setlist{itemsep=2pt, topsep=3pt}

\pagestyle{fancy}
\fancyhf{}
\lhead{\small BTS CIEL 1 -- Gestion de projet}
\rhead{\small TP Git \& GitHub\ifcorrige{} -- \textcolor{red!70!black}{corrigé}\fi}
\cfoot{\small \thepage/\pageref{LastPage}}

% Commande en ligne : \detokenize permet d'écrire _ ~ : sans échappement
% (pas d'accents ni de # dans l'argument).
\newcommand{\cmd}[1]{\texttt{\textcolor{ciel!70!black}{\detokenize{#1}}}}
% Élément de l'interface web de GitHub
\newcommand{\gh}[1]{\textsf{\textbf{#1}}}

\lstdefinestyle{bash}{
  basicstyle=\ttfamily\small, backgroundcolor=\color{gray!8},
  frame=single, rulecolor=\color{gray!40}, framesep=4pt,
  xleftmargin=4pt, xrightmargin=4pt,
  columns=fullflexible, keepspaces=true, breaklines=true,
  showstringspaces=false, upquote=true, aboveskip=6pt, belowskip=6pt,
  literate={é}{{\'e}}1 {è}{{\`e}}1 {ê}{{\^e}}1 {à}{{\`a}}1 {â}{{\^a}}1
           {ù}{{\`u}}1 {û}{{\^u}}1 {ô}{{\^o}}1 {î}{{\^i}}1 {ç}{{\c{c}}}1
           {É}{{\'E}}1 {È}{{\`E}}1 {À}{{\`A}}1
}
\lstdefinestyle{sortie}{style=bash, backgroundcolor=\color{white},
  frame=leftline, framerule=2pt, rulecolor=\color{ciel!50}, basicstyle=\ttfamily\small\color{black!75}}
\lstset{style=bash}

\newtcolorbox{info}[1][À retenir]{breakable, colback=ciel!5, colframe=ciel!80!black,
  fonttitle=\bfseries, title=#1}
\newtcolorbox{attention}[1][Attention]{breakable, colback=red!4, colframe=red!60!black,
  fonttitle=\bfseries, title=#1}
\newtcolorbox[auto counter]{exercice}[1]{breakable, colback=white, colframe=orange!80!black,
  fonttitle=\bfseries, title={Exercice~\thetcbcounter~--~#1}}

\tikzset{
  zone/.style={draw, thick, rounded corners, minimum width=2.9cm, minimum height=1.4cm,
               align=center, fill=#1, font=\small},
  >={Stealth[length=2.5mm]}
}

\begin{document}

% =====================================================================
\begin{tcolorbox}[colback=ciel!6, colframe=ciel, arc=2mm]
  {\LARGE\bfseries TP -- Premiers pas avec Git et GitHub}\\[1mm]
  {\large Sauvegarder et partager son travail, la bonne routine}\\[2mm]
  BTS CIEL -- 1\textsuperscript{re} année \hfill Gestion de projet
  \ifcorrige\\[2mm]{\bfseries\color{red!70!black} VERSION CORRIGÉE -- DOCUMENT ENSEIGNANT}\fi
  \tcblower
  \begin{tabularx}{\linewidth}{@{}>{\bfseries}l X@{}}
    Durée      & 2 h à 3 h \\
    Prérequis  & ouvrir un terminal, utiliser \cmd{cd} et \cmd{ls}, éditer un fichier avec \cmd{nano} ou VS Code \\
    Matériel   & poste avec Git (Git Bash sous Windows), navigateur web, adresse e-mail pour créer un compte GitHub \\
  \end{tabularx}
\end{tcolorbox}

\ifcorrige\else
\vspace{2mm}
\textbf{Nom :} \makebox[6cm]{\dotfill} \hfill \textbf{Prénom :} \makebox[6cm]{\dotfill}
\vspace{2mm}
\fi

\textbf{Objectif.} Ce TP ne fait pas de vous des experts de Git. Il vous apprend \textbf{une routine} que vous
utiliserez dans tous vos TP et dans votre projet de 2\textsuperscript{e} année :
\begin{center}
\cmd{git pull} $\rightarrow$ travailler $\rightarrow$ \cmd{git status} $\rightarrow$ \cmd{git add}
$\rightarrow$ \cmd{git commit} $\rightarrow$ \cmd{git push}
\end{center}
À la fin du TP, votre travail sera sauvegardé sur GitHub, avec l'historique de chaque modification.

% =====================================================================
\section{Git et GitHub en deux mots}

\begin{itemize}
  \item \textbf{Git} est un logiciel installé sur votre PC. Il enregistre l'historique de vos fichiers : chaque
  sauvegarde s'appelle un \textbf{commit} (une « photo » de vos fichiers, avec un message et votre nom).
  \item \textbf{GitHub} est un site web qui héberge une copie de votre travail. Vous le retrouvez depuis n'importe
  quel PC, et vous pouvez le partager avec votre équipe.
  \item Un \textbf{dépôt} (\emph{repository}, ou \emph{repo}) est un dossier de projet suivi par Git. Il existe en
  deux exemplaires : le dépôt \textbf{distant} sur GitHub et le dépôt \textbf{local} sur votre PC.
\end{itemize}

\begin{center}
\begin{tikzpicture}[node distance=1.55cm]
  \node[zone=blue!8]                  (wd)  {Fichiers\\modifiés};
  \node[zone=orange!12, right=of wd]  (idx) {Prêts à être\\enregistrés};
  \node[zone=yellow!15, right=of idx] (loc) {Historique\\sur le PC};
  \node[zone=green!12, right=of loc]  (gh)  {GitHub\\(en ligne)};
  \draw[->, thick] (wd.north east)  to[bend left=35] node[above, font=\small]{\texttt{git add}}    (idx.north west);
  \draw[->, thick] (idx.north east) to[bend left=35] node[above, font=\small]{\texttt{git commit}} (loc.north west);
  \draw[->, thick] (loc.north east) to[bend left=35] node[above, font=\small]{\texttt{git push}}   (gh.north west);
  \draw[->, thick, dashed] (gh.south) to[bend left=14] node[below=1pt, font=\small]{\texttt{git pull} : récupérer les nouveautés de GitHub} (wd.south);
  \draw[gray, rounded corners, densely dotted] ([xshift=-2mm,yshift=11mm]wd.north west) rectangle ([xshift=2mm,yshift=-3mm]loc.south east);
  \node[font=\footnotesize\itshape, text=gray!60!black, anchor=north west] at ([xshift=-1mm,yshift=10mm]wd.north west) {sur votre PC};
\end{tikzpicture}
\end{center}

\begin{info}[Commit ou push ?]
\cmd{git commit} enregistre une version \emph{sur votre PC}. Tant que vous n'avez pas fait \cmd{git push},
GitHub ne la connaît pas : si le PC tombe en panne ou si vous changez de poste, elle est perdue.
\end{info}

% =====================================================================
\section{Préparer GitHub et son poste}

\begin{exercice}{Créer son compte et son premier dépôt sur GitHub}
\begin{enumerate}
  \item Sur \url{https://github.com}, créez un compte (\gh{Sign up}). Choisissez un nom d'utilisateur sérieux :
  il pourra apparaître sur votre CV. Notez-le. \reponse{1}
  \item Créez un dépôt : bouton \gh{+} en haut à droite, puis \gh{New repository}.
  \begin{itemize}
    \item \emph{Repository name} : \cmd{tp-git} ;
    \item visibilité : \gh{Private} ;
    \item activez l'option qui ajoute un fichier \gh{README} ;
    \item validez avec \gh{Create repository}.
  \end{itemize}
  \item Que contient votre dépôt ? Cliquez sur l'horloge « \gh{1 Commit} » : quel est le message de ce commit ? \reponse{2}
\end{enumerate}
\begin{correction}
\begin{enumerate}
  \item[3.] Un seul fichier, \cmd{README.md}, créé par GitHub avec le commit \emph{Initial commit}.
  Le dépôt a donc déjà un historique d'un commit.
\end{enumerate}
\end{correction}
\end{exercice}

\textbf{Configurer Git sur le poste.} Git inscrit votre nom et votre e-mail dans chaque commit. Ouvrez
\textbf{Git Bash} (Windows) ou un terminal (Linux) et tapez, avec \emph{vos} informations :
\begin{lstlisting}
git --version
git config --global user.name "Prénom Nom"
git config --global user.email "adresse-du-compte-github@exemple.fr"
git config --global core.editor nano
git config --global pull.rebase false
\end{lstlisting}
Les deux dernières lignes règlent le comportement de Git quand il fusionne votre travail avec celui d'un camarade
(partie 6). Vérifiez avec \cmd{git config --global --list}.

\begin{attention}[Poste partagé]
Sur un PC de la salle utilisé par d'autres élèves, vérifiez en début de séance que
\cmd{git config user.name} affiche bien \emph{votre} nom. Sinon, refaites la configuration.
\end{attention}

% =====================================================================
\section{Récupérer son dépôt sur le PC : git clone}

Sur la page GitHub de votre dépôt, cliquez sur le bouton vert \gh{Code}, onglet \gh{HTTPS}, et copiez l'adresse.
Dans le terminal (dans Git Bash, collez avec \textbf{Maj+Inser} ou clic droit $\rightarrow$ \emph{Paste}) :
\begin{lstlisting}
cd ~
git clone https://github.com/VOTRE-COMPTE/tp-git.git
cd tp-git
ls
git status
\end{lstlisting}

\cmd{git clone} ne se fait \textbf{qu'une seule fois} par poste : il crée le dossier \cmd{tp-git}, relié à GitHub.

\begin{info}[Connexion à GitHub]
La première fois, Git demande de vous connecter. Sous Windows, une fenêtre s'ouvre : choisissez
\gh{Sign in with your browser} et autorisez l'accès. Le mot de passe du compte n'est pas accepté directement
dans le terminal. Sous Linux, utilisez un \emph{jeton d'accès personnel} (\emph{personal access token}) créé dans
\gh{Settings} $\rightarrow$ \gh{Developer settings}.
\end{info}

\begin{exercice}{Cloner son dépôt}
\begin{enumerate}
  \item Clonez votre dépôt avec les commandes ci-dessus. Quel(s) fichier(s) voyez-vous avec \cmd{ls} ? \reponse{1}
  \item Recopiez les deux premières lignes affichées par \cmd{git status}. \reponse{2}
\end{enumerate}
\begin{correction}
\begin{enumerate}
  \item \cmd{README.md}. Il existe aussi un dossier caché \cmd{.git} (visible avec \cmd{ls -a}) : c'est là que Git
  range l'historique. On n'y touche jamais.
  \item \cmd{On branch main} puis \cmd{Your branch is up to date with 'origin/main'.} (en français :
  « Sur la branche main. Votre branche est à jour avec 'origin/main'. »). \cmd{origin} désigne GitHub.
  En dernière ligne : \emph{nothing to commit, working tree clean}, rien à enregistrer.
\end{enumerate}
\end{correction}
\end{exercice}

% =====================================================================
\section{La routine : status, add, commit, push}

\begin{tabularx}{\linewidth}{@{}>{\raggedright\arraybackslash}p{5.3cm} X@{}}
\toprule
\textbf{Commande} & \textbf{Rôle} \\
\midrule
\cmd{git status}               & « Où en suis-je ? » : quels fichiers sont nouveaux, modifiés, prêts à être enregistrés. \textbf{À taper tout le temps.} \\
\cmd{git add fichier}          & sélectionne un fichier pour le prochain commit (\cmd{git add .} sélectionne tout le dossier) \\
\cmd{git commit -m "message"}  & enregistre les fichiers sélectionnés, avec un message qui dit ce qui a changé \\
\cmd{git push}                 & envoie les commits vers GitHub \\
\bottomrule
\end{tabularx}

\begin{exercice}{Premier cycle complet (guidé)}
Vous allez documenter le réseau de la salle B204.
\begin{enumerate}
  \item Dans le dossier \cmd{tp-git}, créez le fichier \cmd{plan_adressage.txt} avec \cmd{nano plan_adressage.txt}
  (enregistrer : \textbf{Ctrl+O}, \textbf{Entrée} ; quitter : \textbf{Ctrl+X}) :
\begin{lstlisting}
Plan d'adressage - salle B204
VLAN 10  ADMIN    192.168.10.0/24  passerelle 192.168.10.254
VLAN 20  ELEVES   192.168.20.0/24  passerelle 192.168.20.254
\end{lstlisting}
  \item \cmd{git status} : dans quelle rubrique, et de quelle couleur, apparaît le fichier ? \reponse{1}
  \item \cmd{git add plan_adressage.txt}, puis \cmd{git status}. Qu'est-ce qui a changé ? \reponse{1}
  \item Enregistrez :
\begin{lstlisting}
git commit -m "Ajout du plan d'adressage de la salle B204"
git status
\end{lstlisting}
  Que vous dit maintenant \cmd{git status} ? Quelle commande vous suggère-t-il ? \reponse{2}
  \item Envoyez vers GitHub avec \cmd{git push}, puis \cmd{git status}. Rafraîchissez la page GitHub : que voyez-vous ?
  \reponse{2}
\end{enumerate}
\begin{correction}
\begin{enumerate}
  \item[2.] \emph{Untracked files} (fichiers non suivis), en rouge : Git voit le fichier mais ne le suit pas encore.
  \item[3.] Le fichier passe dans \emph{Changes to be committed}, en vert : il sera dans le prochain commit.
  \item[4.] \cmd{Your branch is ahead of 'origin/main' by 1 commit.} suivi de
  \cmd{(use "git push" to publish your local commits)}. Le commit existe sur le PC, pas encore sur GitHub.
  \item[5.] \cmd{git status} revient à \emph{up to date with 'origin/main'}. Sur GitHub, le fichier
  \cmd{plan_adressage.txt} apparaît avec le message du commit, et le compteur passe à 2 commits.
\end{enumerate}
\end{correction}
\end{exercice}

\begin{info}[Lire git status]
\begin{tabularx}{\linewidth}{@{}>{\raggedright\arraybackslash}p{7cm} X@{}}
\emph{Untracked files} (rouge)            & nouveau fichier, pas encore suivi $\rightarrow$ \cmd{git add} \\
\emph{Changes not staged for commit} (rouge) & fichier modifié, pas encore sélectionné $\rightarrow$ \cmd{git add} \\
\emph{Changes to be committed} (vert)     & prêt à être enregistré $\rightarrow$ \cmd{git commit} \\
\emph{Your branch is ahead of 'origin/main'} & commits pas encore envoyés $\rightarrow$ \cmd{git push} \\
\emph{nothing to commit, working tree clean} & tout est enregistré \\
\end{tabularx}
\end{info}

\begin{exercice}{À vous de jouer (sans aide)}
\begin{enumerate}
  \item Ajoutez au plan d'adressage la ligne suivante :
\begin{lstlisting}
VLAN 30  SERVEURS 192.168.30.0/24  passerelle 192.168.30.254
\end{lstlisting}
  \item Créez un fichier \cmd{inventaire.csv} :
\begin{lstlisting}
nom;type;ip
SW-B204;switch;192.168.10.2
PC-B204-01;poste;192.168.20.11
\end{lstlisting}
  \item Faites \textbf{deux commits séparés} (un par modification), avec des messages clairs, puis envoyez-les sur
  GitHub. Notez ci-dessous toutes les commandes utilisées, dans l'ordre. \reponse{7}
  \item Sur GitHub, cliquez sur le compteur de commits. Combien y en a-t-il ? \reponse{1}
\end{enumerate}
\begin{correction}
\begin{lstlisting}
nano plan_adressage.txt
git status
git add plan_adressage.txt
git commit -m "Ajout du VLAN 30 (serveurs)"
nano inventaire.csv
git add inventaire.csv
git commit -m "Ajout de l'inventaire du matériel"
git push
\end{lstlisting}
Un seul \cmd{git push} suffit pour envoyer les deux commits. GitHub affiche 4 commits : \emph{Initial commit} et les
trois commits de l'élève. On valorise les \cmd{git status} intercalés.
\end{correction}
\end{exercice}

\textbf{Voir l'historique.} Sur le PC, \cmd{git log --oneline} affiche un commit par ligne, du plus récent au plus
ancien. Sur GitHub, le compteur de commits donne la même liste. En cliquant sur un commit, on voit les lignes
ajoutées (en vert) et supprimées (en rouge).

\begin{exercice}{Bon ou mauvais message ?}
Un bon message dit \emph{ce qui a changé}, en une phrase courte. Pour chaque message, indiquez B (bon) ou M (mauvais).

\renewcommand{\arraystretch}{1.5}
\begin{tabularx}{\linewidth}{@{}l X c@{}}
\toprule
& \textbf{Message} & \textbf{B/M} \\
\midrule
1 & « modif » & \ifcorrige M\fi \\
2 & « Ajout du VLAN 30 (serveurs) au plan d'adressage » & \ifcorrige B\fi \\
3 & « azerty » & \ifcorrige M\fi \\
4 & « Correction de l'adresse IP du switch dans l'inventaire » & \ifcorrige B\fi \\
5 & « travail de mardi » & \ifcorrige M\fi \\
\bottomrule
\end{tabularx}
\renewcommand{\arraystretch}{1}
\begin{correction}
Les messages 1, 3 et 5 ne disent rien de la modification : dans trois mois, personne (ni vous) ne saura ce qu'ils
contiennent.
\end{correction}
\end{exercice}

% =====================================================================
\section{Récupérer les nouveautés : git pull}

Quand le dépôt GitHub a changé (modification par un camarade, ou par vous depuis un autre PC), on récupère ces
changements avec \cmd{git pull}.

\begin{exercice}{Modifier sur GitHub, récupérer sur le PC}
\begin{enumerate}
  \item Sur GitHub, ouvrez \cmd{README.md}, cliquez sur le crayon (\gh{Edit this file}) et ajoutez la ligne :
  \emph{Documentation du réseau de la salle B204 -- Prénom Nom}. Validez avec \gh{Commit changes}.
  \item Sur le PC, tapez \cmd{git status}. Git signale-t-il la modification faite sur GitHub ? \reponse{1}
  \item Tapez \cmd{git pull}, puis affichez le README avec \cmd{cat README.md}. \reponse{1}
  \item Pourquoi \cmd{git status} ne voyait-il pas la modification à la question 2 ? \reponse{2}
\end{enumerate}
\begin{correction}
\begin{enumerate}
  \item[2.] Non : il affiche toujours \emph{up to date with 'origin/main'}.
  \item[3.] Git affiche \emph{Fast-forward} et le README contient la nouvelle ligne.
  \item[4.] \cmd{git status} ne contacte pas GitHub : il compare avec le dernier état de GitHub \emph{connu} du PC.
  Seuls \cmd{git pull} (et \cmd{git fetch}) vont chercher les nouveautés. D'où la règle :
  \textbf{toujours commencer par \cmd{git pull}}.
\end{enumerate}
\end{correction}
\end{exercice}

\begin{attention}[Jamais de mot de passe sur GitHub]
Tout ce qui est poussé sur GitHub reste dans l'historique, même si l'on supprime le fichier ensuite.
Ne versionnez jamais un mot de passe, une clé ou une configuration contenant un secret. Si cela arrive,
le mot de passe doit être considéré comme compromis et changé.
\end{attention}

% =====================================================================
\section{Travailler à deux sur le même dépôt}

\begin{exercice}{En binôme (élèves A et B)}
\begin{enumerate}
  \item \textbf{(A)} Sur la page GitHub de votre dépôt : \gh{Settings} $\rightarrow$ \gh{Collaborators}
  $\rightarrow$ \gh{Add people}, et ajoutez le compte de B. \textbf{(B)} Acceptez l'invitation reçue par e-mail
  (ou sur \url{https://github.com/notifications}).
  \item \textbf{(B)} Clonez le dépôt de A dans un nouveau dossier :
\begin{lstlisting}
cd ~
git clone https://github.com/COMPTE-DE-A/tp-git.git tp-git-A
cd tp-git-A
\end{lstlisting}
  \item \textbf{(B)} Ajoutez la ligne \cmd{SW-B204-2;switch;192.168.10.3} à la fin de \cmd{inventaire.csv}, puis
  faites la routine jusqu'au \cmd{git push}.
  \item \textbf{(A)} Faites \cmd{git pull} et vérifiez que la ligne de B est bien dans votre inventaire. \reponse{1}
  \item \textbf{En même temps} : A ajoute une ligne au README, B ajoute un VLAN au plan d'adressage. Chacun fait
  \cmd{add} et \cmd{commit}. B fait \cmd{git push} en premier, \emph{puis} A. Que se passe-t-il pour A ? \reponse{2}
  \item \textbf{(A)} Faites \cmd{git pull}. Si \cmd{nano} s'ouvre avec un message de fusion, enregistrez
  (\textbf{Ctrl+O}, \textbf{Entrée}) et quittez (\textbf{Ctrl+X}). Puis \cmd{git push}. Tout le monde fait ensuite
  \cmd{git pull} : les deux dépôts sont-ils identiques ? \reponse{1}
\end{enumerate}
\begin{correction}
\begin{enumerate}
  \item[4.] La ligne de B apparaît chez A.
  \item[5.] Le \cmd{git push} de A est refusé : \cmd{! [rejected] main -> main (fetch first)}. GitHub contient un
  commit (celui de B) que A n'a pas encore. Rien n'est perdu.
  \item[6.] \cmd{git pull} récupère le commit de B et le fusionne avec celui de A (fichiers différents : fusion
  automatique). Après le \cmd{git push} de A et le \cmd{git pull} de B, les deux dépôts sont identiques.
\end{enumerate}
Si les deux élèves modifient \emph{la même ligne}, Git signale un \emph{CONFLICT} : hors programme de ce TP,
l'enseignant aide à le résoudre (garder la bonne ligne, supprimer les marqueurs \cmd{<<<<<<<}, \cmd{=======},
\cmd{>>>>>>>}, puis \cmd{add}, \cmd{commit}, \cmd{push}).
\end{correction}
\end{exercice}

\begin{info}[« push rejected » : pas de panique]
Si \cmd{git push} est refusé, c'est que GitHub a des nouveautés que vous n'avez pas. Faites \cmd{git pull}, puis
\cmd{git push}. Si Git parle de \emph{CONFLICT}, appelez l'enseignant.
\end{info}

% =====================================================================
\section{Défi final (en autonomie)}

Sans regarder les pages précédentes, en vous aidant seulement de l'aide-mémoire :
\begin{enumerate}
  \item Créez sur GitHub un nouveau dépôt privé \cmd{mes-tp-ciel}, avec un README, et clonez-le.
  \item Ajoutez un fichier \cmd{compte_rendu.md} contenant le titre de ce TP et trois choses que vous avez apprises.
  Commit et push.
  \item Complétez le README avec votre nom et votre classe. Commit et push.
  \item Modifiez \cmd{compte_rendu.md} directement sur GitHub, puis récupérez la modification sur le PC.
  \item Montrez à l'enseignant la liste des commits sur GitHub et un \cmd{git status} « propre ».
\end{enumerate}
Ce dépôt vous servira pour ranger vos prochains TP.

\begin{correction}
Attendus : au moins 3 commits (\emph{Initial commit}, \cmd{compte_rendu.md}, README) plus le commit fait sur GitHub,
visible sur le PC après \cmd{git pull}. \cmd{git status} doit afficher \emph{up to date with 'origin/main'} et
\emph{nothing to commit, working tree clean}. Vérifier aussi la qualité des messages.
\end{correction}

\subsection*{Auto-évaluation}
\begin{tabularx}{\linewidth}{@{}X c c c@{}}
\toprule
\textbf{Je sais\ldots} & \textbf{Pas encore} & \textbf{En partie} & \textbf{Oui} \\
\midrule
expliquer la différence entre Git et GitHub, entre commit et push & $\square$ & $\square$ & $\square$ \\
créer un dépôt sur GitHub et le cloner sur un PC & $\square$ & $\square$ & $\square$ \\
lire \cmd{git status} et savoir quelle commande taper ensuite & $\square$ & $\square$ & $\square$ \\
enregistrer mon travail : \cmd{add}, \cmd{commit} avec un message clair & $\square$ & $\square$ & $\square$ \\
envoyer et récupérer mon travail : \cmd{push}, \cmd{pull} & $\square$ & $\square$ & $\square$ \\
réagir à un \emph{push rejected} & $\square$ & $\square$ & $\square$ \\
\bottomrule
\end{tabularx}

% =====================================================================
\vspace{5mm}
\needspace{20\baselineskip}
\section*{Aide-mémoire}
\addcontentsline{toc}{section}{Aide-mémoire}

\begin{tcolorbox}[colback=ciel!5, colframe=ciel, title=La routine de chaque séance, fonttitle=\bfseries]
\begin{tabularx}{\linewidth}{@{}l >{\raggedright\arraybackslash}p{5.5cm} X@{}}
\textbf{1.} & \cmd{git pull}                 & récupérer les nouveautés (début de séance) \\
\textbf{2.} & \emph{travailler}              & créer, modifier des fichiers \\
\textbf{3.} & \cmd{git status}               & voir ce qui a changé \\
\textbf{4.} & \cmd{git add fichier}          & sélectionner (ou \cmd{git add .} pour tout) \\
\textbf{5.} & \cmd{git commit -m "message"}  & enregistrer sur le PC \\
\textbf{6.} & \cmd{git push}                 & envoyer sur GitHub (au plus tard en fin de séance) \\
\end{tabularx}
\end{tcolorbox}

\begin{tabularx}{\linewidth}{@{}>{\raggedright\arraybackslash}p{7.3cm} X@{}}
\toprule
\multicolumn{2}{@{}l}{\textbf{\color{ciel}Une seule fois}} \\
\cmd{git config --global user.name "Nom"} & votre nom (dans chaque commit) \\
\cmd{git config --global user.email "mail"} & votre e-mail (celui du compte GitHub) \\
\cmd{git clone <url>} & copier un dépôt GitHub sur le PC \\
\midrule
\multicolumn{2}{@{}l}{\textbf{\color{ciel}Pour s'y retrouver}} \\
\cmd{git status} & où en suis-je ? \\
\cmd{git log --oneline} & liste des commits \\
\bottomrule
\end{tabularx}

\subsection*{Pour aller plus loin}
\begin{itemize}
  \item S. Chacon et B. Straub, \emph{Pro Git}, 2\textsuperscript{e} éd., Apress, 2014. Libre, traduit en français :
  \url{https://git-scm.com/book/fr/v2} (chapitres 1 et 2).
  \item Documentation GitHub en français : \url{https://docs.github.com/fr/get-started}.
\end{itemize}

\end{document}
"
%  Compiler deux fois (numéro de dernière page dans le pied de page).
% =====================================================================
\documentclass[11pt,a4paper]{article}

\usepackage[utf8]{inputenc}
\usepackage[T1]{fontenc}
\usepackage[french]{babel}
\usepackage{lmodern}
\usepackage{microtype}
% Pas de ligatures -- et << >> dans les commandes (police à chasse fixe)
\DisableLigatures[-,<,>]{encoding=T1, family=tt*}
\usepackage{needspace}
\usepackage[margin=2.2cm,headheight=14pt]{geometry}
\usepackage{xcolor}
\usepackage{upquote}
\usepackage{listings}
\usepackage[most]{tcolorbox}
\usepackage{tikz}
\usetikzlibrary{positioning,arrows.meta,babel}
\usepackage{enumitem}
\usepackage{tabularx}
\usepackage{booktabs}
\usepackage{amssymb}
\usepackage{fancyhdr}
\usepackage{lastpage}
\usepackage{comment}
\usepackage[hidelinks]{hyperref}

% ---------------------------------------------------------------------
% Version élève / version corrigée
% ---------------------------------------------------------------------
\newif\ifcorrige
\ifdefined\CORRIGE\corrigetrue\fi

\definecolor{ciel}{RGB}{0,84,147}

\ifcorrige
  \newtcolorbox{correction}{breakable, colback=green!4, colframe=green!45!black,
    fonttitle=\bfseries\small, title=Correction, before skip=2mm, after skip=3mm}
  \newcommand{\reponse}[1]{}
\else
  \excludecomment{correction}
  \newcommand{\reponse}[1]{\par\smallskip{\color{gray!70}%
    \foreach \n in {1,...,#1}{\noindent\mbox{}\dotfill\par\vspace{2.5mm}}}}
\fi

% ---------------------------------------------------------------------
% Mise en forme
% ---------------------------------------------------------------------
\setlength{\parindent}{0pt}
\setlength{\parskip}{4pt}
\setlist{itemsep=2pt, topsep=3pt}

\pagestyle{fancy}
\fancyhf{}
\lhead{\small BTS CIEL 1 -- Gestion de projet}
\rhead{\small TP Git \& GitHub\ifcorrige{} -- \textcolor{red!70!black}{corrigé}\fi}
\cfoot{\small \thepage/\pageref{LastPage}}

% Commande en ligne : \detokenize permet d'écrire _ ~ : sans échappement
% (pas d'accents ni de # dans l'argument).
\newcommand{\cmd}[1]{\texttt{\textcolor{ciel!70!black}{\detokenize{#1}}}}
% Élément de l'interface web de GitHub
\newcommand{\gh}[1]{\textsf{\textbf{#1}}}

\lstdefinestyle{bash}{
  basicstyle=\ttfamily\small, backgroundcolor=\color{gray!8},
  frame=single, rulecolor=\color{gray!40}, framesep=4pt,
  xleftmargin=4pt, xrightmargin=4pt,
  columns=fullflexible, keepspaces=true, breaklines=true,
  showstringspaces=false, upquote=true, aboveskip=6pt, belowskip=6pt,
  literate={é}{{\'e}}1 {è}{{\`e}}1 {ê}{{\^e}}1 {à}{{\`a}}1 {â}{{\^a}}1
           {ù}{{\`u}}1 {û}{{\^u}}1 {ô}{{\^o}}1 {î}{{\^i}}1 {ç}{{\c{c}}}1
           {É}{{\'E}}1 {È}{{\`E}}1 {À}{{\`A}}1
}
\lstdefinestyle{sortie}{style=bash, backgroundcolor=\color{white},
  frame=leftline, framerule=2pt, rulecolor=\color{ciel!50}, basicstyle=\ttfamily\small\color{black!75}}
\lstset{style=bash}

\newtcolorbox{info}[1][À retenir]{breakable, colback=ciel!5, colframe=ciel!80!black,
  fonttitle=\bfseries, title=#1}
\newtcolorbox{attention}[1][Attention]{breakable, colback=red!4, colframe=red!60!black,
  fonttitle=\bfseries, title=#1}
\newtcolorbox[auto counter]{exercice}[1]{breakable, colback=white, colframe=orange!80!black,
  fonttitle=\bfseries, title={Exercice~\thetcbcounter~--~#1}}

\tikzset{
  zone/.style={draw, thick, rounded corners, minimum width=2.9cm, minimum height=1.4cm,
               align=center, fill=#1, font=\small},
  >={Stealth[length=2.5mm]}
}

\begin{document}

% =====================================================================
\begin{tcolorbox}[colback=ciel!6, colframe=ciel, arc=2mm]
  {\LARGE\bfseries TP -- Premiers pas avec Git et GitHub}\\[1mm]
  {\large Sauvegarder et partager son travail, la bonne routine}\\[2mm]
  BTS CIEL -- 1\textsuperscript{re} année \hfill Gestion de projet
  \ifcorrige\\[2mm]{\bfseries\color{red!70!black} VERSION CORRIGÉE -- DOCUMENT ENSEIGNANT}\fi
  \tcblower
  \begin{tabularx}{\linewidth}{@{}>{\bfseries}l X@{}}
    Durée      & 2 h à 3 h \\
    Prérequis  & ouvrir un terminal, utiliser \cmd{cd} et \cmd{ls}, éditer un fichier avec \cmd{nano} ou VS Code \\
    Matériel   & poste avec Git (Git Bash sous Windows), navigateur web, adresse e-mail pour créer un compte GitHub \\
  \end{tabularx}
\end{tcolorbox}

\ifcorrige\else
\vspace{2mm}
\textbf{Nom :} \makebox[6cm]{\dotfill} \hfill \textbf{Prénom :} \makebox[6cm]{\dotfill}
\vspace{2mm}
\fi

\textbf{Objectif.} Ce TP ne fait pas de vous des experts de Git. Il vous apprend \textbf{une routine} que vous
utiliserez dans tous vos TP et dans votre projet de 2\textsuperscript{e} année :
\begin{center}
\cmd{git pull} $\rightarrow$ travailler $\rightarrow$ \cmd{git status} $\rightarrow$ \cmd{git add}
$\rightarrow$ \cmd{git commit} $\rightarrow$ \cmd{git push}
\end{center}
À la fin du TP, votre travail sera sauvegardé sur GitHub, avec l'historique de chaque modification.

% =====================================================================
\section{Git et GitHub en deux mots}

\begin{itemize}
  \item \textbf{Git} est un logiciel installé sur votre PC. Il enregistre l'historique de vos fichiers : chaque
  sauvegarde s'appelle un \textbf{commit} (une « photo » de vos fichiers, avec un message et votre nom).
  \item \textbf{GitHub} est un site web qui héberge une copie de votre travail. Vous le retrouvez depuis n'importe
  quel PC, et vous pouvez le partager avec votre équipe.
  \item Un \textbf{dépôt} (\emph{repository}, ou \emph{repo}) est un dossier de projet suivi par Git. Il existe en
  deux exemplaires : le dépôt \textbf{distant} sur GitHub et le dépôt \textbf{local} sur votre PC.
\end{itemize}

\begin{center}
\begin{tikzpicture}[node distance=1.55cm]
  \node[zone=blue!8]                  (wd)  {Fichiers\\modifiés};
  \node[zone=orange!12, right=of wd]  (idx) {Prêts à être\\enregistrés};
  \node[zone=yellow!15, right=of idx] (loc) {Historique\\sur le PC};
  \node[zone=green!12, right=of loc]  (gh)  {GitHub\\(en ligne)};
  \draw[->, thick] (wd.north east)  to[bend left=35] node[above, font=\small]{\texttt{git add}}    (idx.north west);
  \draw[->, thick] (idx.north east) to[bend left=35] node[above, font=\small]{\texttt{git commit}} (loc.north west);
  \draw[->, thick] (loc.north east) to[bend left=35] node[above, font=\small]{\texttt{git push}}   (gh.north west);
  \draw[->, thick, dashed] (gh.south) to[bend left=14] node[below=1pt, font=\small]{\texttt{git pull} : récupérer les nouveautés de GitHub} (wd.south);
  \draw[gray, rounded corners, densely dotted] ([xshift=-2mm,yshift=11mm]wd.north west) rectangle ([xshift=2mm,yshift=-3mm]loc.south east);
  \node[font=\footnotesize\itshape, text=gray!60!black, anchor=north west] at ([xshift=-1mm,yshift=10mm]wd.north west) {sur votre PC};
\end{tikzpicture}
\end{center}

\begin{info}[Commit ou push ?]
\cmd{git commit} enregistre une version \emph{sur votre PC}. Tant que vous n'avez pas fait \cmd{git push},
GitHub ne la connaît pas : si le PC tombe en panne ou si vous changez de poste, elle est perdue.
\end{info}

% =====================================================================
\section{Préparer GitHub et son poste}

\begin{exercice}{Créer son compte et son premier dépôt sur GitHub}
\begin{enumerate}
  \item Sur \url{https://github.com}, créez un compte (\gh{Sign up}). Choisissez un nom d'utilisateur sérieux :
  il pourra apparaître sur votre CV. Notez-le. \reponse{1}
  \item Créez un dépôt : bouton \gh{+} en haut à droite, puis \gh{New repository}.
  \begin{itemize}
    \item \emph{Repository name} : \cmd{tp-git} ;
    \item visibilité : \gh{Private} ;
    \item activez l'option qui ajoute un fichier \gh{README} ;
    \item validez avec \gh{Create repository}.
  \end{itemize}
  \item Que contient votre dépôt ? Cliquez sur l'horloge « \gh{1 Commit} » : quel est le message de ce commit ? \reponse{2}
\end{enumerate}
\begin{correction}
\begin{enumerate}
  \item[3.] Un seul fichier, \cmd{README.md}, créé par GitHub avec le commit \emph{Initial commit}.
  Le dépôt a donc déjà un historique d'un commit.
\end{enumerate}
\end{correction}
\end{exercice}

\textbf{Configurer Git sur le poste.} Git inscrit votre nom et votre e-mail dans chaque commit. Ouvrez
\textbf{Git Bash} (Windows) ou un terminal (Linux) et tapez, avec \emph{vos} informations :
\begin{lstlisting}
git --version
git config --global user.name "Prénom Nom"
git config --global user.email "adresse-du-compte-github@exemple.fr"
git config --global core.editor nano
git config --global pull.rebase false
\end{lstlisting}
Les deux dernières lignes règlent le comportement de Git quand il fusionne votre travail avec celui d'un camarade
(partie 6). Vérifiez avec \cmd{git config --global --list}.

\begin{attention}[Poste partagé]
Sur un PC de la salle utilisé par d'autres élèves, vérifiez en début de séance que
\cmd{git config user.name} affiche bien \emph{votre} nom. Sinon, refaites la configuration.
\end{attention}

% =====================================================================
\section{Récupérer son dépôt sur le PC : git clone}

Sur la page GitHub de votre dépôt, cliquez sur le bouton vert \gh{Code}, onglet \gh{HTTPS}, et copiez l'adresse.
Dans le terminal (dans Git Bash, collez avec \textbf{Maj+Inser} ou clic droit $\rightarrow$ \emph{Paste}) :
\begin{lstlisting}
cd ~
git clone https://github.com/VOTRE-COMPTE/tp-git.git
cd tp-git
ls
git status
\end{lstlisting}

\cmd{git clone} ne se fait \textbf{qu'une seule fois} par poste : il crée le dossier \cmd{tp-git}, relié à GitHub.

\begin{info}[Connexion à GitHub]
La première fois, Git demande de vous connecter. Sous Windows, une fenêtre s'ouvre : choisissez
\gh{Sign in with your browser} et autorisez l'accès. Le mot de passe du compte n'est pas accepté directement
dans le terminal. Sous Linux, utilisez un \emph{jeton d'accès personnel} (\emph{personal access token}) créé dans
\gh{Settings} $\rightarrow$ \gh{Developer settings}.
\end{info}

\begin{exercice}{Cloner son dépôt}
\begin{enumerate}
  \item Clonez votre dépôt avec les commandes ci-dessus. Quel(s) fichier(s) voyez-vous avec \cmd{ls} ? \reponse{1}
  \item Recopiez les deux premières lignes affichées par \cmd{git status}. \reponse{2}
\end{enumerate}
\begin{correction}
\begin{enumerate}
  \item \cmd{README.md}. Il existe aussi un dossier caché \cmd{.git} (visible avec \cmd{ls -a}) : c'est là que Git
  range l'historique. On n'y touche jamais.
  \item \cmd{On branch main} puis \cmd{Your branch is up to date with 'origin/main'.} (en français :
  « Sur la branche main. Votre branche est à jour avec 'origin/main'. »). \cmd{origin} désigne GitHub.
  En dernière ligne : \emph{nothing to commit, working tree clean}, rien à enregistrer.
\end{enumerate}
\end{correction}
\end{exercice}

% =====================================================================
\section{La routine : status, add, commit, push}

\begin{tabularx}{\linewidth}{@{}>{\raggedright\arraybackslash}p{5.3cm} X@{}}
\toprule
\textbf{Commande} & \textbf{Rôle} \\
\midrule
\cmd{git status}               & « Où en suis-je ? » : quels fichiers sont nouveaux, modifiés, prêts à être enregistrés. \textbf{À taper tout le temps.} \\
\cmd{git add fichier}          & sélectionne un fichier pour le prochain commit (\cmd{git add .} sélectionne tout le dossier) \\
\cmd{git commit -m "message"}  & enregistre les fichiers sélectionnés, avec un message qui dit ce qui a changé \\
\cmd{git push}                 & envoie les commits vers GitHub \\
\bottomrule
\end{tabularx}

\begin{exercice}{Premier cycle complet (guidé)}
Vous allez documenter le réseau de la salle B204.
\begin{enumerate}
  \item Dans le dossier \cmd{tp-git}, créez le fichier \cmd{plan_adressage.txt} avec \cmd{nano plan_adressage.txt}
  (enregistrer : \textbf{Ctrl+O}, \textbf{Entrée} ; quitter : \textbf{Ctrl+X}) :
\begin{lstlisting}
Plan d'adressage - salle B204
VLAN 10  ADMIN    192.168.10.0/24  passerelle 192.168.10.254
VLAN 20  ELEVES   192.168.20.0/24  passerelle 192.168.20.254
\end{lstlisting}
  \item \cmd{git status} : dans quelle rubrique, et de quelle couleur, apparaît le fichier ? \reponse{1}
  \item \cmd{git add plan_adressage.txt}, puis \cmd{git status}. Qu'est-ce qui a changé ? \reponse{1}
  \item Enregistrez :
\begin{lstlisting}
git commit -m "Ajout du plan d'adressage de la salle B204"
git status
\end{lstlisting}
  Que vous dit maintenant \cmd{git status} ? Quelle commande vous suggère-t-il ? \reponse{2}
  \item Envoyez vers GitHub avec \cmd{git push}, puis \cmd{git status}. Rafraîchissez la page GitHub : que voyez-vous ?
  \reponse{2}
\end{enumerate}
\begin{correction}
\begin{enumerate}
  \item[2.] \emph{Untracked files} (fichiers non suivis), en rouge : Git voit le fichier mais ne le suit pas encore.
  \item[3.] Le fichier passe dans \emph{Changes to be committed}, en vert : il sera dans le prochain commit.
  \item[4.] \cmd{Your branch is ahead of 'origin/main' by 1 commit.} suivi de
  \cmd{(use "git push" to publish your local commits)}. Le commit existe sur le PC, pas encore sur GitHub.
  \item[5.] \cmd{git status} revient à \emph{up to date with 'origin/main'}. Sur GitHub, le fichier
  \cmd{plan_adressage.txt} apparaît avec le message du commit, et le compteur passe à 2 commits.
\end{enumerate}
\end{correction}
\end{exercice}

\begin{info}[Lire git status]
\begin{tabularx}{\linewidth}{@{}>{\raggedright\arraybackslash}p{7cm} X@{}}
\emph{Untracked files} (rouge)            & nouveau fichier, pas encore suivi $\rightarrow$ \cmd{git add} \\
\emph{Changes not staged for commit} (rouge) & fichier modifié, pas encore sélectionné $\rightarrow$ \cmd{git add} \\
\emph{Changes to be committed} (vert)     & prêt à être enregistré $\rightarrow$ \cmd{git commit} \\
\emph{Your branch is ahead of 'origin/main'} & commits pas encore envoyés $\rightarrow$ \cmd{git push} \\
\emph{nothing to commit, working tree clean} & tout est enregistré \\
\end{tabularx}
\end{info}

\begin{exercice}{À vous de jouer (sans aide)}
\begin{enumerate}
  \item Ajoutez au plan d'adressage la ligne suivante :
\begin{lstlisting}
VLAN 30  SERVEURS 192.168.30.0/24  passerelle 192.168.30.254
\end{lstlisting}
  \item Créez un fichier \cmd{inventaire.csv} :
\begin{lstlisting}
nom;type;ip
SW-B204;switch;192.168.10.2
PC-B204-01;poste;192.168.20.11
\end{lstlisting}
  \item Faites \textbf{deux commits séparés} (un par modification), avec des messages clairs, puis envoyez-les sur
  GitHub. Notez ci-dessous toutes les commandes utilisées, dans l'ordre. \reponse{7}
  \item Sur GitHub, cliquez sur le compteur de commits. Combien y en a-t-il ? \reponse{1}
\end{enumerate}
\begin{correction}
\begin{lstlisting}
nano plan_adressage.txt
git status
git add plan_adressage.txt
git commit -m "Ajout du VLAN 30 (serveurs)"
nano inventaire.csv
git add inventaire.csv
git commit -m "Ajout de l'inventaire du matériel"
git push
\end{lstlisting}
Un seul \cmd{git push} suffit pour envoyer les deux commits. GitHub affiche 4 commits : \emph{Initial commit} et les
trois commits de l'élève. On valorise les \cmd{git status} intercalés.
\end{correction}
\end{exercice}

\textbf{Voir l'historique.} Sur le PC, \cmd{git log --oneline} affiche un commit par ligne, du plus récent au plus
ancien. Sur GitHub, le compteur de commits donne la même liste. En cliquant sur un commit, on voit les lignes
ajoutées (en vert) et supprimées (en rouge).

\begin{exercice}{Bon ou mauvais message ?}
Un bon message dit \emph{ce qui a changé}, en une phrase courte. Pour chaque message, indiquez B (bon) ou M (mauvais).

\renewcommand{\arraystretch}{1.5}
\begin{tabularx}{\linewidth}{@{}l X c@{}}
\toprule
& \textbf{Message} & \textbf{B/M} \\
\midrule
1 & « modif » & \ifcorrige M\fi \\
2 & « Ajout du VLAN 30 (serveurs) au plan d'adressage » & \ifcorrige B\fi \\
3 & « azerty » & \ifcorrige M\fi \\
4 & « Correction de l'adresse IP du switch dans l'inventaire » & \ifcorrige B\fi \\
5 & « travail de mardi » & \ifcorrige M\fi \\
\bottomrule
\end{tabularx}
\renewcommand{\arraystretch}{1}
\begin{correction}
Les messages 1, 3 et 5 ne disent rien de la modification : dans trois mois, personne (ni vous) ne saura ce qu'ils
contiennent.
\end{correction}
\end{exercice}

% =====================================================================
\section{Récupérer les nouveautés : git pull}

Quand le dépôt GitHub a changé (modification par un camarade, ou par vous depuis un autre PC), on récupère ces
changements avec \cmd{git pull}.

\begin{exercice}{Modifier sur GitHub, récupérer sur le PC}
\begin{enumerate}
  \item Sur GitHub, ouvrez \cmd{README.md}, cliquez sur le crayon (\gh{Edit this file}) et ajoutez la ligne :
  \emph{Documentation du réseau de la salle B204 -- Prénom Nom}. Validez avec \gh{Commit changes}.
  \item Sur le PC, tapez \cmd{git status}. Git signale-t-il la modification faite sur GitHub ? \reponse{1}
  \item Tapez \cmd{git pull}, puis affichez le README avec \cmd{cat README.md}. \reponse{1}
  \item Pourquoi \cmd{git status} ne voyait-il pas la modification à la question 2 ? \reponse{2}
\end{enumerate}
\begin{correction}
\begin{enumerate}
  \item[2.] Non : il affiche toujours \emph{up to date with 'origin/main'}.
  \item[3.] Git affiche \emph{Fast-forward} et le README contient la nouvelle ligne.
  \item[4.] \cmd{git status} ne contacte pas GitHub : il compare avec le dernier état de GitHub \emph{connu} du PC.
  Seuls \cmd{git pull} (et \cmd{git fetch}) vont chercher les nouveautés. D'où la règle :
  \textbf{toujours commencer par \cmd{git pull}}.
\end{enumerate}
\end{correction}
\end{exercice}

\begin{attention}[Jamais de mot de passe sur GitHub]
Tout ce qui est poussé sur GitHub reste dans l'historique, même si l'on supprime le fichier ensuite.
Ne versionnez jamais un mot de passe, une clé ou une configuration contenant un secret. Si cela arrive,
le mot de passe doit être considéré comme compromis et changé.
\end{attention}

% =====================================================================
\section{Travailler à deux sur le même dépôt}

\begin{exercice}{En binôme (élèves A et B)}
\begin{enumerate}
  \item \textbf{(A)} Sur la page GitHub de votre dépôt : \gh{Settings} $\rightarrow$ \gh{Collaborators}
  $\rightarrow$ \gh{Add people}, et ajoutez le compte de B. \textbf{(B)} Acceptez l'invitation reçue par e-mail
  (ou sur \url{https://github.com/notifications}).
  \item \textbf{(B)} Clonez le dépôt de A dans un nouveau dossier :
\begin{lstlisting}
cd ~
git clone https://github.com/COMPTE-DE-A/tp-git.git tp-git-A
cd tp-git-A
\end{lstlisting}
  \item \textbf{(B)} Ajoutez la ligne \cmd{SW-B204-2;switch;192.168.10.3} à la fin de \cmd{inventaire.csv}, puis
  faites la routine jusqu'au \cmd{git push}.
  \item \textbf{(A)} Faites \cmd{git pull} et vérifiez que la ligne de B est bien dans votre inventaire. \reponse{1}
  \item \textbf{En même temps} : A ajoute une ligne au README, B ajoute un VLAN au plan d'adressage. Chacun fait
  \cmd{add} et \cmd{commit}. B fait \cmd{git push} en premier, \emph{puis} A. Que se passe-t-il pour A ? \reponse{2}
  \item \textbf{(A)} Faites \cmd{git pull}. Si \cmd{nano} s'ouvre avec un message de fusion, enregistrez
  (\textbf{Ctrl+O}, \textbf{Entrée}) et quittez (\textbf{Ctrl+X}). Puis \cmd{git push}. Tout le monde fait ensuite
  \cmd{git pull} : les deux dépôts sont-ils identiques ? \reponse{1}
\end{enumerate}
\begin{correction}
\begin{enumerate}
  \item[4.] La ligne de B apparaît chez A.
  \item[5.] Le \cmd{git push} de A est refusé : \cmd{! [rejected] main -> main (fetch first)}. GitHub contient un
  commit (celui de B) que A n'a pas encore. Rien n'est perdu.
  \item[6.] \cmd{git pull} récupère le commit de B et le fusionne avec celui de A (fichiers différents : fusion
  automatique). Après le \cmd{git push} de A et le \cmd{git pull} de B, les deux dépôts sont identiques.
\end{enumerate}
Si les deux élèves modifient \emph{la même ligne}, Git signale un \emph{CONFLICT} : hors programme de ce TP,
l'enseignant aide à le résoudre (garder la bonne ligne, supprimer les marqueurs \cmd{<<<<<<<}, \cmd{=======},
\cmd{>>>>>>>}, puis \cmd{add}, \cmd{commit}, \cmd{push}).
\end{correction}
\end{exercice}

\begin{info}[« push rejected » : pas de panique]
Si \cmd{git push} est refusé, c'est que GitHub a des nouveautés que vous n'avez pas. Faites \cmd{git pull}, puis
\cmd{git push}. Si Git parle de \emph{CONFLICT}, appelez l'enseignant.
\end{info}

% =====================================================================
\section{Défi final (en autonomie)}

Sans regarder les pages précédentes, en vous aidant seulement de l'aide-mémoire :
\begin{enumerate}
  \item Créez sur GitHub un nouveau dépôt privé \cmd{mes-tp-ciel}, avec un README, et clonez-le.
  \item Ajoutez un fichier \cmd{compte_rendu.md} contenant le titre de ce TP et trois choses que vous avez apprises.
  Commit et push.
  \item Complétez le README avec votre nom et votre classe. Commit et push.
  \item Modifiez \cmd{compte_rendu.md} directement sur GitHub, puis récupérez la modification sur le PC.
  \item Montrez à l'enseignant la liste des commits sur GitHub et un \cmd{git status} « propre ».
\end{enumerate}
Ce dépôt vous servira pour ranger vos prochains TP.

\begin{correction}
Attendus : au moins 3 commits (\emph{Initial commit}, \cmd{compte_rendu.md}, README) plus le commit fait sur GitHub,
visible sur le PC après \cmd{git pull}. \cmd{git status} doit afficher \emph{up to date with 'origin/main'} et
\emph{nothing to commit, working tree clean}. Vérifier aussi la qualité des messages.
\end{correction}

\subsection*{Auto-évaluation}
\begin{tabularx}{\linewidth}{@{}X c c c@{}}
\toprule
\textbf{Je sais\ldots} & \textbf{Pas encore} & \textbf{En partie} & \textbf{Oui} \\
\midrule
expliquer la différence entre Git et GitHub, entre commit et push & $\square$ & $\square$ & $\square$ \\
créer un dépôt sur GitHub et le cloner sur un PC & $\square$ & $\square$ & $\square$ \\
lire \cmd{git status} et savoir quelle commande taper ensuite & $\square$ & $\square$ & $\square$ \\
enregistrer mon travail : \cmd{add}, \cmd{commit} avec un message clair & $\square$ & $\square$ & $\square$ \\
envoyer et récupérer mon travail : \cmd{push}, \cmd{pull} & $\square$ & $\square$ & $\square$ \\
réagir à un \emph{push rejected} & $\square$ & $\square$ & $\square$ \\
\bottomrule
\end{tabularx}

% =====================================================================
\vspace{5mm}
\needspace{20\baselineskip}
\section*{Aide-mémoire}
\addcontentsline{toc}{section}{Aide-mémoire}

\begin{tcolorbox}[colback=ciel!5, colframe=ciel, title=La routine de chaque séance, fonttitle=\bfseries]
\begin{tabularx}{\linewidth}{@{}l >{\raggedright\arraybackslash}p{5.5cm} X@{}}
\textbf{1.} & \cmd{git pull}                 & récupérer les nouveautés (début de séance) \\
\textbf{2.} & \emph{travailler}              & créer, modifier des fichiers \\
\textbf{3.} & \cmd{git status}               & voir ce qui a changé \\
\textbf{4.} & \cmd{git add fichier}          & sélectionner (ou \cmd{git add .} pour tout) \\
\textbf{5.} & \cmd{git commit -m "message"}  & enregistrer sur le PC \\
\textbf{6.} & \cmd{git push}                 & envoyer sur GitHub (au plus tard en fin de séance) \\
\end{tabularx}
\end{tcolorbox}

\begin{tabularx}{\linewidth}{@{}>{\raggedright\arraybackslash}p{7.3cm} X@{}}
\toprule
\multicolumn{2}{@{}l}{\textbf{\color{ciel}Une seule fois}} \\
\cmd{git config --global user.name "Nom"} & votre nom (dans chaque commit) \\
\cmd{git config --global user.email "mail"} & votre e-mail (celui du compte GitHub) \\
\cmd{git clone <url>} & copier un dépôt GitHub sur le PC \\
\midrule
\multicolumn{2}{@{}l}{\textbf{\color{ciel}Pour s'y retrouver}} \\
\cmd{git status} & où en suis-je ? \\
\cmd{git log --oneline} & liste des commits \\
\bottomrule
\end{tabularx}

\subsection*{Pour aller plus loin}
\begin{itemize}
  \item S. Chacon et B. Straub, \emph{Pro Git}, 2\textsuperscript{e} éd., Apress, 2014. Libre, traduit en français :
  \url{https://git-scm.com/book/fr/v2} (chapitres 1 et 2).
  \item Documentation GitHub en français : \url{https://docs.github.com/fr/get-started}.
\end{itemize}

\end{document}
